\documentclass[11pt]{article}
\usepackage[margin=1in]{geometry}
\usepackage{fontspec}
\usepackage{titlesec}
\usepackage{enumitem}
\usepackage{longtable}
\usepackage{array}
\usepackage{booktabs}
\usepackage{xcolor}
\usepackage{hyperref}
\usepackage{fancyhdr}
\usepackage{tcolorbox}
\usepackage{amssymb}
\usepackage{amsmath}
\usepackage{siunitx}
\usepackage[normalem]{ulem}
\usepackage{listings}
\usepackage{graphicx}

\lstset{
  basicstyle=\fontsize{6pt}{7pt}\ttfamily,
  breaklines=true,
  breakatwhitespace=false,
  columns=fullflexible,
  frame=single,
  framerule=0.3pt,
  rulecolor=\color{gray!40},
}
\tcbuselibrary{skins,breakable}

\definecolor{marlinblue}{HTML}{1B4965}
\definecolor{marlinteal}{HTML}{2E8B8B}
\definecolor{lightgray}{HTML}{F4F4F4}
\definecolor{codebg}{HTML}{F7F7F9}
\definecolor{risk}{HTML}{B23A48}

\titleformat{\section}{\large\bfseries\color{marlinblue}}{\thesection}{0.5em}{}[{\color{marlinblue}\titlerule}]
\titleformat{\subsection}{\normalsize\bfseries\color{marlinteal}}{\thesubsection}{0.5em}{}

\hypersetup{colorlinks=true, linkcolor=marlinblue, urlcolor=marlinteal}

\pagestyle{fancy}
\fancyhf{}
\fancyhead[L]{\small\textcolor{gray}{Marlin FD + PETSc Roadmap}}
\fancyhead[R]{\small\textcolor{gray}{\thepage}}
\renewcommand{\headrulewidth}{0.3pt}

\setlength{\parskip}{0.5em}
\setlength{\parindent}{0em}
\sloppy
\emergencystretch=3em
\Urlmuskip=0mu plus 1mu

\newenvironment{schematic}{%
  \begin{tcolorbox}[breakable, colback=codebg, colframe=gray!40, boxrule=0.4pt, arc=1mm, left=4mm, right=4mm, top=2mm, bottom=2mm]
  \ttfamily\footnotesize
}{%
  \end{tcolorbox}
}

\newcommand{\pr}[2]{\subsection{PR#1 --- #2}}
\newcommand{\status}[1]{\textcolor{marlinteal}{\textbf{Status:}} #1}

\title{\textcolor{marlinblue}{\Huge Marlin Finite-Difference + PETSc\\[2pt] Implicit Solve}\\[8pt]
\Large Development Roadmap}
\author{}
\date{\today}

\begin{document}
\maketitle

\begin{tcolorbox}[colback=lightgray, colframe=marlinblue, boxrule=0.6pt, arc=1mm]
\textbf{Status of this document: record and remaining plan.} It records the
decisions taken, what was built and where, and a short plan for the PRs that
remain. The code is the source of truth: where this document and
\texttt{src/}, \texttt{include/} or the tests disagree, the code wins and this
document is corrected. Each PR carries a status (\texttt{proposed} /
\texttt{approved} / \texttt{implemented} / \texttt{verified} / \texttt{merged}).
Detailed design is kept only for PRs that remain; built PRs carry a short
description of what was built.
\end{tcolorbox}

\tableofcontents
\vspace{1em}

\section{Context and Goal}

Marlin composes PDE physics as \texttt{TensorCompute} objects
(\texttt{include/tensor\_computes/TensorOperatorBase.h}). Before this work the
real-space path had one FD operator (\texttt{FiniteDifferenceLaplacian})
stepped explicitly (\texttt{RealSpaceForwardEuler}), no PETSc in the tensor
solve path, and no boundary-condition mechanism: smooth boundary method (SBM)
no-flux is imposed by a multiplicative \texttt{psi} mask on the right-hand
side (\texttt{test/tests/kks/KKS\_no\_flux\_bc.i}).

\textbf{Goal:} compose PDE terms as TensorComputes, recognise the terms that
are linear with a known stencil, and solve the implied $Ax=b$ implicitly with
a Krylov method whose operator action reuses the explicit kernels
(matrix-free, Section~\ref{sec:partial-assembly}); lag the other terms into
$b$; Picard-iterate where coupling needs it; allow steady fields (e.g.\ a
Poisson equation) alongside transient ones (Section~\ref{sec:steady-mixing}).
Target functionality: Cahn-Hilliard spinodal decomposition, KKS /
grand-potential, and the smooth boundary method with real no-flux boundary
conditions.

\textbf{Built so far} (Section~\ref{sec:roadmap}, Status Tracking): the
\texttt{LinearFDOperator} interface, real-space Dirichlet / Neumann / Robin
boundary conditions, \texttt{TimeDerivativeOperator},
\texttt{FiniteDifferenceBiharmonic}, and \texttt{ImplicitFDSolver}
(\texttt{src/tensor\_solver/ImplicitFDSolver.C}) with a PETSc KSPCG + PCMG
backend (CPU, CUDA) and a torch PCG backend (any device, default on MPS), on
periodic and non-periodic domains. Remaining: variable-coefficient diffusion,
coupled steady + transient fields, Picard coupling, no-flux biharmonic, and
the KKS/SBM port.

The original plan came from a survey of the TensorCompute / TensorSolver /
TensorProblem / DomainAction architecture, MOOSE's \texttt{PetscSupport}
options mechanism and \texttt{FixedPointSolve} Picard knobs, the target
example inputs (\texttt{test/tests/kks/KKS\_no\_flux\_bc.i},
\texttt{examples/cahn\_hilliard/*.i} including \texttt{cahnhilliard\_fem.i},
\texttt{examples/libtorch\_kks/KKS\_libtorch.i}), decisions with the user
(Section~\ref{sec:decisions}), and a risk review (Section~\ref{sec:risks}).

\section{Architecture as Built}
\label{sec:architecture}

This section describes the code on branch \texttt{implicit\_fd\_bcs}
(PR0--PR4). Parts that are planned are labelled as such.

\subsection{Hybrid Design: PETSc Owns the Algorithms, Torch Owns the Kernels}
\label{sec:hybrid}

\begin{tcolorbox}[colback=lightgray, colframe=marlinblue, boxrule=0.6pt, arc=1mm]
\textbf{Decided 2026-10-18} after the PR1a/PR1b evidence chain
(Section~\ref{sec:pr1b}); hierarchy ownership revised after PR2b (user
decision, same day). The deciding factor was generality and long-term
maintenance (PETSc's solver stack is maintained by full-time scientists),
not speed: the measured gap was not large in either direction.
\end{tcolorbox}

As built in \texttt{ImplicitFDSolver} (\texttt{src/tensor\_solver/ImplicitFDSolver.C})
and \texttt{PetscFD} (\texttt{src/utils/PetscFDHierarchy.C}), PETSc backend:

\begin{longtable}{p{0.3\textwidth} p{0.1\textwidth} p{0.5\textwidth}}
\toprule
Layer & Owner & Mechanism \\
\midrule
\endhead
Krylov method, convergence test & PETSc & \texttt{KSPCG}, unpreconditioned residual norm, \texttt{l\_tol}, nonzero initial guess (warm start from $u^n$); no divergence test, since $\|r_0\|/\|b\|$ grows with $dt$ (PR3b) \\
Level hierarchy and partitions & PETSc & Fine \texttt{DMDA} on Marlin's partition (ownership checked), coarse DMDAs by \texttt{DMCoarsenHierarchy}; \texttt{DMDA\_Q0} keeps non-periodic axes cell-centred ($n \to n/2$). Depth from \texttt{PetscFD::countLevels}; \texttt{-pc\_mg\_levels} is an error \\
Fine-level operator and diagonal & torch & \texttt{MATSHELL}: \texttt{MatMult} = \texttt{applyLinearOperator} (each operator's own \texttt{computeBuffer()}, ghosts by Marlin's \texttt{exchangeGhostLayers}); \texttt{GetDiagonal} = \texttt{jacobiDiagonal} \\
Coarse-level operators & PETSc & Rediscretised from \texttt{stencilAt(spacing)}, assembled by \texttt{MatSetValuesStencil}, scaled by $2^{\dim k}$ (restriction is $P^T$), preallocated from the stencil's columns; reassembled only when the stencil changes (e.g.\ a new $dt$) \\
Transfers & PETSc (our matrix) & Cell-centred multilinear interpolation, assembled AIJ with global coarse columns (\texttt{PetscFD::cellCentredInterpolation}); restriction $P^T$ \\
Smoothers & PETSc & Default \texttt{smoother = CHEBYSHEV}: PCMG's Chebyshev, SOR on the CPU coarse levels, Jacobi on the fine level and on CUDA, bounds re-estimated on every level when the operators change. \texttt{smoother = DAMPED\_JACOBI}: \texttt{KSPRICHARDSON} + \texttt{PCJACOBI}, weight $\omega = 4/(3\rho)$, $\rho$ the Gershgorin bound of $D^{-1}A$ over all rows (PR3b, PR4) \\
Coarse solve & PETSc & PCMG's default LU; shifted (\texttt{MAT\_SHIFT\_NONZERO}) for a singular field \\
Boundary ghost values & Marlin & \texttt{TensorRobinBC} ghost fill on the fine level; homogeneous ghost factor folded into coarse rows and the interpolation (PR4) \\
Solution / RHS \texttt{Vec}s & PETSc on torch memory & \texttt{VecCreateMPIWithArray} / \texttt{VecCreateMPICUDAWithArray} + \texttt{VecPlaceArray} on the solver's tensors (zero-copy) \\
\bottomrule
\end{longtable}

Configuration rules from the evidence, as applied in the code:
\begin{itemize}[leftmargin=*]
\item \textbf{SPD marking.} Fine and coarse levels of a definite field get
\texttt{MAT\_SPD} and \texttt{MAT\_SPD\_ETERNAL} (\texttt{MatAssemblyEnd}
clears \texttt{MAT\_SPD} otherwise; the coarse flag is set again after every
assembly). A singular field (constants as nullspace) gets
\texttt{MAT\_SYMMETRIC} and a \texttt{MatNullSpace} on every level instead.
In PR1b, marking SPD halved PCMG setup on the H200 (724$\to$408~ms at
$n=8192$) with the same iterations and solution.
\item \textbf{Smoother.} PCMG's Chebyshev is the default
(\texttt{smoother = CHEBYSHEV}), with SOR on the CPU coarse levels and Jacobi on
the fine level and on CUDA, where PETSc's SOR runs on the host. The
\texttt{KSP\_DIVERGED\_INDEFINITE\_PC} seen on CUDA for 1a came from stale bounds:
without a DM, PCMG sets up only the fine smoother again on a new matrix, so the
coarse levels kept the bounds of the first $dt$. Each coarse smoother's operators
are reset after reassembly. \texttt{smoother = DAMPED\_JACOBI} is
Richardson + Jacobi with the Gershgorin weight; a user-set
\texttt{-<field>\_mg\_levels\_ksp\_richardson\_scale} wins. The analytic-bounds
hook (\texttt{KSPChebyshevSetEigenvalues}) was not built.
\item \textbf{Options.} \texttt{petsc\_options*} of the solver block go into
the global options database under the prefix \texttt{<field>\_}; command-line
options win. Fields whose names differ only in case are an error.
\end{itemize}

Torch backend (\texttt{linear\_solver = TORCH\_PCG}, default on MPS):
\texttt{MooseTensor::conjugateGradientSolve} with the same
\texttt{applyLinearOperator} and a Jacobi preconditioner
(\texttt{jacobiDiagonal}). It also runs in single precision; the PETSc
backend requires double precision. The torch smoothers of the PR1b/PR2c spikes are
not in the production path.

\subsection{Operator Application: Matrix-Free on the Fine Level}
\label{sec:partial-assembly}

\begin{tcolorbox}[colback=lightgray, colframe=marlinblue, boxrule=0.6pt, arc=1mm]
\textbf{Decision:} the Krylov operator $A$ is applied matrix-free (MFEM-style
partial assembly), not assembled as in MOOSE. Only the small coarse
multigrid levels are assembled.
\end{tcolorbox}

The fine-level action is exactly the explicit kernel: \texttt{applyOperator}
writes $x$ into the field's owned cells, runs each operator's
\texttt{computeBuffer()} with ghosts (a \texttt{convNd} for the FD
operators), and sums the outputs. Reasons, recorded at decision time: Marlin
already had the matrix-free kernels on a structured grid; it avoids a
round-trip through PETSc's \texttt{Mat} layout on GPU-resident tensors; PETSc
has no Metal/MPS backend, so on MPS only a torch action is possible; the
action works on the whole per-field tensor (Section~\ref{sec:soa}). An
assembled preconditioning matrix for AMG is the contingent PR4a.

\subsection{Data Layout and Zero-Copy at the PETSc Boundary}
\label{sec:soa}

Every \texttt{TensorCompute} buffer is one dense grid-shaped
\texttt{torch::Tensor} per field (structure-of-arrays), unlike MOOSE's
interleaved per-node DOFs. Each field is its own single-DOF PETSc system, so
this layout is kept. A future monolithic solver (PR8) must use field-split /
block \texttt{Vec}s rather than interleaved DOFs.

\textbf{Zero-copy is a requirement} where a field's tensor becomes a PETSc
\texttt{Vec}. As built, the $b$ and $x$ \texttt{Vec}s are created without
storage and \texttt{VecPlaceArray} puts the solver's contiguous tensors in
before every solve (CPU: \texttt{VecCreateMPIWithArray}; CUDA:
\texttt{VecCreateMPICUDAWithArray}). The shell callbacks wrap PETSc's arrays
as tensors with \texttt{torch::from\_blob} on the array's memory type. On
H200 at $n=256$, \texttt{KSPSolve} makes no host/device copies with Jacobi
smoothing on every level (PR3). MPS has no PETSc \texttt{Vec} type, so it
uses the torch backend. Zero-copy memory sharing is not a claim about
matching compute kernels: PETSc's own vector work inside \texttt{KSPSolve}
runs through cuBLAS.

\subsection{Class Schematic}

\begin{schematic}
\begin{verbatim}
TensorOperatorBase
  |
  +-- TensorOperator<T> + LinearFDOperator (mixin: stencilAt(),
  |     |                  stencilAt(spacing), definiteness(),
  |     |                  inputBufferName(), rhsContribution())
  |     +-- FiniteDifferenceLaplacian      (PR0; 3/5-point, factor)
  |     +-- TimeDerivativeOperator         (PR3; u/dt, rhs u^n/dt)
  |     +-- FiniteDifferenceBiharmonic     (PR3b; non-periodic 3-point
  |     |                                   in PR5)
  |     +-- FiniteDifferenceDivGrad        (PR5; factor div(D grad u))
  |
  +-- TensorBoundaryCondition  (variable, boundary face;
  |     |                       fillGhostCells(), prepareFaceData())
  |     +-- TensorRobinBC      (a u + b du/dn = c; a, b scalar; c constant,
  |           |                 c_function or c_buffer; ghostCoefficientOnU0,
  |           |                 ghostConstantTerm)
  |           +-- TensorDirichletBC, TensorNeumannBC (presets)
  |
  +-- TensorSolver
        +-- ImplicitFDSolver   (PR3/PR4; fields 'buffer', per-field
        |                       'operators' and 'source'; backends
        |                       PETSC and TORCH_PCG)
        +-- [planned] Picard iteration over ImplicitFDSolver (PR6)

PetscFD (utils/PetscFDHierarchy): Partition, FaceGhosts, Hierarchy,
  buildHierarchy, assembleStencilOperator, createStencilMatrix,
  jacobiSpectrumBound, cellCentredInterpolation

TensorProblem: [TensorComputes/BCs] via AddTensorBCAction ->
  addTensorComputeBC, getBoundaryCondition(buffer, face),
  exchangeGhostLayers (halo + BC ghost fill), beginFixedBoundaryData /
  endFixedBoundaryData
\end{verbatim}
\end{schematic}

Names in the original design that were not built: \texttt{PetscFDSolverBase}
(built as \texttt{ImplicitFDSolver}); \texttt{applyAt},
\texttt{rhsContributionAt} (built as the operators' own
\texttt{computeBuffer()} and \texttt{rhsContribution()});
\texttt{coefficientsChanged} (the solver compares the summed stencil with the
one last assembled); \texttt{StencilEntry::target\_field} (not built; the
struct is \texttt{FDStencilEntry\{offset, coefficient\}});
\texttt{assembleRHS} (inline in \texttt{substep()});
\texttt{getBoundaryConditionsFor} (built as
\texttt{getBoundaryCondition(buffer, face)});
\texttt{applyGhostContribution} (dropped: the affine split below needs no
separate boundary method); \texttt{stencilContribution} (built in PR1,
removed in PR4 because nothing used it); \texttt{getGlobalIndexOffset}
(not needed: the partition is gathered from public \texttt{DomainAction}
getters). Still planned: \texttt{PicardTensorSolver} (PR6),
\texttt{assemblePreconditionerMatrix}
(PR4a).

\subsection{Data Flow of One Time Step}

\begin{schematic}
\begin{verbatim}
ImplicitFDSolver::substep()
  run the explicit computes (everything in the root compute that is
    not a field operator): sources f and buffer-valued BC data, at t^n
  beginFixedBoundaryData(dt): exchange BC buffer halos once; BC
    functions are evaluated at t^n + dt
  for each field, in the order of 'buffer':
    b = sum_i r_i (e.g. u^n/dt) + f          (u still holds u^n)
    with boundary faces: g = A(0) (ghost fill with the BC values),
                         b -= g, solve the linear part A(x) - A(0)
    singular field: check sum(b) is round-off, remove the mean
    solve A x = b:  PETSC: KSPCG + PCMG, fine shell = applyLinearOperator
                    TORCH_PCG: conjugateGradientSolve, Jacobi
    failed solve -> restore every field, _is_converged = false
                    (TensorSolveIterationAdaptiveDT cuts dt), or error
    write x into u, apply the operators once more (outputs hold A_i u),
    exchangeGhostLayers(u)
  endFixedBoundaryData()
[planned, PR6] Picard: repeat the field loop until converged
\end{verbatim}
\end{schematic}

Fields in one \texttt{ImplicitFDSolver} are solved one after another, but
their sources are computed once, before the loop, from the state at $t^n$.

\subsection{Boundary-Condition Dual-Use Invariant}

This is the design detail most likely to silently rot if not actively
enforced by a test --- flagged explicitly as Risk R4 in
Section~\ref{sec:risks}.

\begin{schematic}
\begin{verbatim}
TensorRobinBC(a, b, c) for buffer "u" on face "left"
ghost = ghostCoefficientOnU0(a, b, s) * mirror + c / (a/2 + b/s),
s = (2l+1) dx for ghost layer l
     |
     +--> fillGhostCells(u)              [ghost fill]
     |      consumed by: explicit computes and solvers, and the
     |                   implicit fine level (matrix-free A(x) - A(0),
     |                   b - A(0): the affine split)
     |
     +--> ghostCoefficientOnU0(a, b, s)  [implicit matrix entries]
            consumed by: the torch PCG Jacobi diagonal, and
                         PetscFD::FaceGhosts::factor on the coarse
                         levels (folded rows, interpolation)

INVARIANT: both derive from one a,b triple and one ghost relation.
REQUIRED TEST (not optional): the ghost fill's dependence on each
mirror cell (fill(u) - fill(0)) equals the coefficients the implicit
solver writes, for both ghost layers and both faces
(TensorBoundaryConditionTest), and the fill satisfies a u + b du/dn = c.
\end{verbatim}
\end{schematic}

\subsection{Steady Fields Mixed with Transient Ones}
\label{sec:steady-mixing}

Requirement (user): a steady PDE (e.g.\ a Poisson equation for an electric
or chemical potential) must be solvable alongside transient phase-field
variables, each feeding source terms into the other.

As built, the time derivative is not an assumption of the solver: it is one
optional \texttt{LinearFDOperator}, \texttt{TimeDerivativeOperator} (action
$u/dt$, \texttt{rhsContribution} $u^n/dt$), listed in a field's
\texttt{operators} like any other term. A steady field lists none; the
periodic Poisson tests (\texttt{test/tests/implicit\_fd/poisson\_mms.i}) solve
such a field, singular with the constants as nullspace. \texttt{ImplicitFDSolver}
accepts several fields (for example \texttt{buffer = 'c phi'} with
\texttt{operators = 'dcdt lap\_c; lap\_phi'}; \texttt{steady\_transient.i}
and \texttt{pfhub\_6a.i}). The explicit computes run once before the field loop, so a
source computed from another field sees that field at $t^n$ even if it was
solved earlier in the same time step: the result does not depend on the order
of \texttt{buffer} (decision D8). Re-solving the steady field against updated
sources is the job of the Picard iteration (PR6); PR5a verifies the coupled
case.

\section{Decisions}
\label{sec:decisions}

Each item records the decision, who made it and when, and its state in the
code.

\begin{enumerate}[leftmargin=*]
\item \textbf{Segregated per-field linear systems in one solver object},
coupled across fields by lagging (and later Picard), not monolithic block
coupling. \emph{Built:} \texttt{ImplicitFDSolver} solves the fields of its
\texttt{buffer} list one after another, each with its own KSP and PETSc
options prefix \texttt{<field>\_}. The \texttt{target\_field} tag on stencil
entries proposed for a future monolithic solver was not built.
\item \textbf{Mutually exclusive solve mode}: a simulation uses one
\texttt{[TensorSolver]}, so implicit FD and the explicit/spectral solvers
are not mixed in one run. \texttt{ImplicitFDSolver} requires
\texttt{Domain/parallel\_mode = REAL\_SPACE}.
\item \textbf{Preconditioner reuse, with an iteration threshold and a
mandatory rebuild-and-retry} (user, 2026-10-18, from the PR1b sweep).
\begin{itemize}
\item \emph{Built:} the shell operator always applies the current
coefficients. The coarse levels and smoother weights are reassembled, and
\texttt{KSPSetOperators} resets every level, only when the field's summed
stencil changes (e.g.\ a new $dt$). A failed solve (negative converged
reason, or a TORCH\_PCG residual above tolerance) restores the fields and
reports non-convergence, so \texttt{TensorSolveIterationAdaptiveDT} retries
with a smaller $dt$; without that time stepper it is an error.
\item \emph{Not built (PR6):} the iteration threshold for rebuilding, and a
rebuild when coefficients that depend on the solution change between Picard
iterations. A rebuild must go through \texttt{KSPSetOperators}: bumping a
shell Mat's state alone re-sets-up only the finest smoother (PETSc 3.25.4
mg.c, independent review).
\item Evidence (PR1b, $\beta=3$, $n=4096$): 1695~ms with reuse vs.\ 2786~ms
with a full rebuild every step. Reuse saves little at $\sim$1300$\times$
coefficient contrast.
\end{itemize}
\item \textbf{Dirichlet, Neumann and Robin} --- Robin is the general
implementation, Dirichlet and Neumann are presets. \emph{Built} for the
explicit path (PR1) and the implicit path (PR4), with the limits listed
under PR4 (5-point Robin is an error).
\item \textbf{Hybrid PETSc + torch} (user, 2026-10-18,
Section~\ref{sec:hybrid}): KSPCG + PCMG with a torch fine-level shell is the
default for SPD fields on CPU and CUDA; PETSc owns the hierarchy, the coarse
levels are rediscretised and assembled, transfers are cell-centred. The
fine operator is matrix-free; an assembled \texttt{Pmat} for AMG is the
contingent PR4a. Zero-copy \texttt{Vec}s are a requirement and are built.
\emph{Built} in PR3.
\item \textbf{Steady fields may be mixed with transient fields} in one
implicit run (Section~\ref{sec:steady-mixing}). This does not reopen item~2.
\item \textbf{MPS through the torch PCG only} (user, 2026-10-18).
\texttt{linear\_solver = AUTO} selects \texttt{TORCH\_PCG} on MPS and
\texttt{PETSC} elsewhere; PETSc requires double precision and, on CUDA, a
CUDA build of PETSc and libtorch. More complex solvers on MPS are deferred.
\item \textbf{PCMG's Chebyshev smoother is the default, damped Jacobi an
option} (user, 2026-10-23, after the \#189 review). On CPU, Chebyshev + SOR kept
the biharmonic iterations flat in $n$ (9 at $n = 64$--$256$, against 15 at
$n = 256$ for damped Jacobi) and took fewer on 120 steps of 1a (mean 4.3 against
5.9). Damped Jacobi with the Gershgorin weight stays as
\texttt{smoother = DAMPED\_JACOBI}. Only the preconditioner changes.
\item \textbf{No PETSc divergence test} (user, 2026-10-24): warm started from
$u^n$, $\|r_0\|/\|b\| \sim dt\,\|Bu^n\|/\|u^n\|$, so a converging large-$dt$ solve
crossed the default \texttt{divtol} at its first iteration. \texttt{TORCH\_PCG}
has none either; \texttt{-<field>\_ksp\_divtol} sets one.
\item \textbf{Coarse multigrid rows next to a boundary face are
rediscretised} (PR4), not Galerkin: ghost layer $l$ is folded onto its
mirror cell with the homogeneous ghost factor
(\texttt{TensorRobinBC::ghostCoefficientOnU0} at spacing $(2l+1)h$ of the
level); the cell-centred interpolation folds the same way and restriction
stays $P^T$; \texttt{DMDA\_Q0} keeps \texttt{DMCoarsen} cell-centred on
non-periodic axes. The boundary values enter only the fine residual. PCMG
iterations with boundary faces 7, 6, 6 at $n = 64, 128, 256$
($l_{tol}=10^{-10}$, CPU) against 9, 9, 9 periodic.
\item \textbf{Boundary value timing} (user, 2026-10-22). In
\texttt{ImplicitFDSolver}, \texttt{Function}-valued boundary values are
evaluated at $t^{n+1}$; buffer-valued ones hold what their compute wrote at
the start of the time step ($t^n$), a one-step lag. This is documented, not
changed: only an explicit function of time can be evaluated at $t^{n+1}$;
removing the other lags is what a Picard iteration does (input to PR6).
Reaffirmed in the PR4 milestone review (user, 2026-10-25), with the test
requirement corrected and \texttt{bcs/be\_3pt\_buffer} checking the lag
against a direct solve with the buffer value at $t^n$ (error 0.0225 against
0.0026 for a function).
\item \textbf{Strict zero-mean check with Neumann faces only} (user,
2026-10-25). Data that balance in the continuous problem generally leave a
discrete mean of order $h^2$; it is an error rather than projected out, so a
genuine inconsistency is not hidden. Documented in \texttt{ImplicitFDSolver.md}.
\item \textbf{A Robin face with $ab < 0$ is an error} (user, 2026-10-25): its
ghost term $a/(a/2 + b/dx)$ turns negative once $dx < 2|b/a|$, after which the
operator may be indefinite (depending on the rest of the problem), so whether a
problem solved depended on the grid; behaviour should not change with $dx$.
\item \textbf{\texttt{FiniteDifferenceDivGrad} face coefficient} (user,
2026-10-25; PR5): harmonic mean $2D_iD_j/(D_i+D_j)$ by default, the exact
series conductance for a jump on a cell face; an arithmetic option for
coefficients that can reach zero (e.g.\ degenerate mobility), where the
harmonic mean shuts the face. The docs say when to use which.
\item \textbf{Sign of $D$} (user, 2026-10-25; PR5): explicit use allows any
sign (negative $D$ is useful e.g.\ as the explicit part of an IMEX split);
harmonic mode requires $D_iD_j > 0$ on every face. In \texttt{ImplicitFDSolver}
$D < 0$ is an error (the operator would not be semidefinite, and whether a
problem solved would depend on $dt$ and $dx$), and so is $D = 0$ in harmonic
mode. Checked at every assembly.
\item \textbf{Biharmonic faces} (user, 2026-10-25; PR5): with the 3-point
Laplacian, Neumann faces impose $\partial u/\partial n = q$ and
$\partial(\nabla^2 u)/\partial n = 0$, Dirichlet faces $u = g$ and
$\nabla^2 u = 0$; Robin faces and the 5-point biharmonic with faces are errors.
The milestone review found that the ghost fill of $u$ alone imposes
$\partial^2u/\partial n^2 = 0$ and $\partial^3u/\partial n^3 = 0$ instead
when the face data vary along the face; the user approved (2026-10-26) setting
the ghosts of $w = \nabla^2 u$ to the exact reflection, which imposes the pairs
above (the homogeneous operator is unchanged), and (2026-10-26) imposing them
wherever the operator is used, explicit and postprocessing included.
\item \textbf{$D$ at a domain face} (user, 2026-10-25; PR5): the boundary
cell's own $D_0$ (a mirrored $D$ ghost, the same under either mean). Exact for
$D$ constant next to the face, $O(h)$ in the face coefficient for a smooth $D$.
A Neumann face still gives $\partial u/\partial n$, so its physical flux is
$D_0$ times the given value; the docs say so. Measured in 1D (NumPy, $n$ = 32 to
256): second order on Dirichlet and no-flux faces, first order on a Neumann face
with a nonzero flux where $D$ varies; kept over linear extrapolation of $D$ (second
order, but it can lose $D$'s sign and reaches across a jump next to the face), and
documented (user, 2026-10-26).
\item \textbf{Coarse $D$} (user, 2026-10-25; PR5): each PCMG coarse level
takes the mean of $D$ over its $2^{\mathrm{dim}}$ child cells and is
rediscretised with the same face rule, as the constant-coefficient levels are.
Preconditioner only. Iterations across jumps are recorded in the capacitor
tests; Galerkin $P^TAP$ or AMG (PR4a) is the follow-up if they degrade.
\item \textbf{Central advection at high cell P\'eclet deferred, CUDA graph
capture not adopted} (user, 2026-10-18; PR2c).
\end{enumerate}

\section{Risks}
\label{sec:risks}

\begin{longtable}{>{\bfseries}p{0.05\textwidth} p{0.3\textwidth} p{0.2\textwidth} p{0.35\textwidth}}
\toprule
\# & Risk & Why it matters & Status \\
\midrule
\endhead
R1 & DMDA partition may not match Marlin's tensor layout rank for rank &
Silent cross-rank corruption & \textbf{Retired.} PR2 spike on 1--8 ranks;
in production \texttt{PetscFD::buildHierarchy} checks the fine DMDA owns
exactly this rank's Marlin block, and \texttt{FineDMDAMatchesMarlinBlock}
tests it \\
R2 & Picard may converge slowly or not at all for strongly nonlinear
coupling (e.g.\ KKS mobility) & Could force a late redesign of the coupling
& \textbf{Open.} PR6, measured on a real coupled case \\
R3 & PCMG setup may dominate Picard runtime; reuse may degrade the
preconditioner at high contrast & ``Performant'' goal & \textbf{Partly
measured} (PR1b: after \texttt{MAT\_SPD} setup is about one solve at
$n\ge4096$; reuse saved 39\% at $\beta=3$). Setup is paid only when the
stencil changes in the current code; re-measure in PR6 \\
R4 & Ghost-fill and implicit-matrix BC math could diverge as code evolves &
Inconsistent boundary treatment, as with the \texttt{psi} mask & \textbf{Retired
for the code as built} (PR4): both use \texttt{ghostCoefficientOnU0};
\texttt{TensorBoundaryConditionTest} checks that the fill's dependence on
each mirror cell equals those coefficients, both layers, both faces. The PR1
dual-path test checked \texttt{stencilContribution}, which nothing used; it
was replaced. Any new operator or BC that writes matrix entries must extend
that test \\
R5 & Robin generality could delay Dirichlet/Neumann & SBM no-flux needs only
Neumann & \textbf{Retired} (PR1) \\
R6 & Torch/PETSc hand-off overhead could erase PETSc's benefit & Premise of
using PETSc on GPU & \textbf{Retired 2026-10-18} (PR1a, PR1b): zero-copy on
CPU and CUDA; at steady state PETSc is not slower (Jacobi-CG 3--19\% faster;
PCMG repeat solve 0.70--0.92$\times$ torch MG-PCG). PR3 confirmed no
host/device copies in \texttt{KSPSolve} on H200 \\
R7 & Geometric MG with point Jacobi may converge slowly on strongly
anisotropic or heterogeneous coefficients & Could force AMG and an
assembled \texttt{Pmat} & \textbf{Open, contingent.} PR1a: AMG needed
1.5$\times$ fewer iterations at $D_x/D_y=10^4$ but was slower at $n=64$.
Revisit in PR5/PR6 with real coefficients; PR4a if needed \\
R8 & PCMG may not extend to multiple ranks (halos, coarse gather, thin
slabs) & Multi-rank is required & \textbf{Retired for the built scope.}
PR2b (PETSc-owned hierarchy); PR3 tests on 1, 2 and 4 ranks; PR4
(non-periodic) multi-rank per the user. Not shown: large rank counts \\
R9 & Non-SPD operators break CG and Jacobi smoothing & SPD default does not
cover all target physics & \textbf{Open.} PR2c spike: Gauss-Seidel robust
for upwind advection; central advection at high Pe deferred.
\texttt{ImplicitFDSolver} accepts only SPD/SPSD operators (an error
otherwise) \\
R10 & Fusion headroom of torch smoothers inside PETSc callbacks & Whether
custom smoothers are worth maintaining & \textbf{Answered} (PR2c): traced
2.4--3.2$\times$ eager on H200, none on CPU. No torch smoother is in the
production path \\
R11 & Lagged boundary data and sources: buffer-valued BC values and explicit
sources are at $t^n$ & Splitting error in coupled or BC-driven problems &
\textbf{Documented} (decision~11); addressed by Picard (PR6) \\
\bottomrule
\end{longtable}

\section{Verification Methodology}
\label{sec:verification}

Every PR that touches numerics must include:

\begin{enumerate}[leftmargin=*]
\item \textbf{Method of Manufactured Solutions (MMS).} Pick a smooth
analytic $u_{\text{exact}}(x,y,t)$, substitute into the governing equation
to derive the forcing/source term, inject that source as an added
\texttt{ParsedCompute}-style term in $b$ (natural fit --- nonlinear/explicit
terms already flow into $b$ this way), run at 2--3 grid resolutions, and
confirm the solution error converges at the expected order (3-pt Laplacian:
$O(dx^2)$; 5-pt: $O(dx^4)$; with boundary faces, record the observed order
rather than assume it).
\item \textbf{Cross-validation against an independent method}, not just
regression against ourselves. Concrete anchors already in the repo:
\begin{itemize}
\item \texttt{examples/cahn\_hilliard/cahnhilliard\_fem.i} --- classic MOOSE
FEM Kernels/BCs implementation of Cahn-Hilliard. Use for PR5/PR7.
\item Marlin's own existing spectral operators
(\texttt{ReciprocalMatDiffusion}, \texttt{ReciprocalAllenCahn}) on a
periodic domain, where both the new FD path and the existing spectral path
should agree within truncation error. Use for PR3/PR5.
\item \texttt{examples/libtorch\_kks/KKS\_libtorch.i} and
\texttt{test/tests/kks/KKS\_no\_flux\_bc.i} (spectral + \texttt{psi}-mask) as
the qualitative/quantitative reference for PR6/PR7's ported KKS/SBM case.
\end{itemize}
\item \textbf{Written record attached to the PR}: convergence-order table or
plot, and the cross-validation comparison (plot or error table), not just
``tests pass.''
\end{enumerate}

\section{Development Process and Gates}
\label{sec:gates}

Every PR in the roadmap goes through the same four-stage gate before merge.
None of these stages are optional, and none can be skipped by a
plausible-sounding justification.

\begin{schematic}
\begin{verbatim}
 Implement --> Independent Reviewer --> Physics Verification --> Merge
                      Agent                     Gate
   (scope        (fresh agent/subagent,    (MMS convergence +     (status
    limited        NO shared context        cross-validation,     updated to
    to this        with implementer,        written record         "verified")
    PR's stated    runs code-review          attached to PR)
    deliverable)   skill adversarially)
\end{verbatim}
\end{schematic}

\begin{description}[leftmargin=*,style=nextline]
\item[Implement] Scope strictly to the PR's stated deliverable. No bundling
of unrelated cleanup. If something consequential is discovered
mid-implementation (e.g.\ R1 turning out false), stop and update this
document's risk/roadmap sections before continuing --- don't silently absorb
it into the PR.
\item[Independent reviewer agent] A fresh agent (not the one that wrote the
code, no inherited conversation context) runs the \texttt{/code-review}
skill on the diff, looking for correctness bugs first,
simplification/efficiency second. Findings must be fixed or explicitly
accepted with a written reason in the PR description --- ``an independent
agent reviewed this and found nothing'' is only true if one actually ran,
fresh, on the real diff.
\item[Physics verification gate] Run the Section~\ref{sec:verification}
procedure for this PR's new numerics; attach results. A PR with new
numerics and no attached convergence/cross-validation evidence does not
pass this gate, regardless of code-review outcome.
\item[Merge / go-no-go] Only after both the review and the verification
gate pass does the PR merge and its status move to \texttt{verified}. A
``go'' on a plan section is not a ``go'' on an unspecified implementation
detail discovered later --- new consequential choices get their own
sign-off.
\end{description}

\textbf{Lessons from PR1--PR4.}
\begin{itemize}[leftmargin=*]
\item \textbf{Tests must check what the solver actually uses.} The PR1 R4
dual-path test checked \texttt{stencilContribution}, a method this roadmap
required but no solver called; it protected an unused path while the
implicit solver used \texttt{ghostCoefficientOnU0}. PR4 removed the method
and tests the coefficients the solver writes. A roadmap requirement for a
method is not a reason to keep it.
\item \textbf{Measure before setting test thresholds.} Iteration-count and
error thresholds come from a measured run on each device, not from an
expected value.
\item \textbf{Every pre-push review finding is confirmed by an experiment}
before it is fixed or dismissed.
\item \textbf{Devices.} CPU and MPS are tested locally; CUDA (H200) runs on
another machine and is tested by the user.
\end{itemize}

\section{PR Roadmap}
\label{sec:roadmap}

Each PR is scoped to be independently reviewable and, where possible,
behavior-neutral for existing functionality until it's explicitly
activated.

\pr{0}{\texttt{LinearFDOperator} Mixin + Retrofit \texttt{FiniteDifferenceLaplacian}}
\textbf{Built:} the header-only mixin
\texttt{include/tensor\_computes/LinearFDOperator.h} (stencil as
\texttt{FDStencilEntry\{offset, coefficient\}} from \texttt{stencilAt()}),
implemented by \texttt{FiniteDifferenceLaplacian} from the same coefficients
as its conv kernel; behavior-neutral. PR3 extended the interface with
\texttt{stencilAt(spacing)}, \texttt{definiteness()},
\texttt{inputBufferName()} and \texttt{rhsContribution()}.

\textbf{Verification:} \texttt{unit/src/LinearFDOperatorTest.C}: stencil
coefficients equal the conv-kernel coefficients for 3- and 5-point stencils
in 1D/2D/3D.

\status{merged into \texttt{devel} (commit \texttt{2580075}, via PR \#185).}

\pr{1}{Boundary Condition TensorComputes (ghost fill, explicit path)}
\textbf{Built:} \texttt{TensorBoundaryCondition},
\texttt{TensorRobinBC} ($a u + b\,\partial u/\partial n = c$, up to two ghost
layers, the second from the pair (\texttt{u1}, ghost 2) at spacing
$3\Delta x$), and the presets \texttt{TensorDirichletBC} /
\texttt{TensorNeumannBC}; the \texttt{[TensorComputes/BCs]} block
(\texttt{AddTensorBCAction}, \texttt{TensorProblem::addTensorComputeBC},
\texttt{getBoundaryCondition(buffer, face)}); \texttt{exchangeGhostLayers}
fills physical-face ghosts independent of the rank count; setup validation
(missing, duplicate or periodic faces). The planned
\texttt{stencilContribution()} was built here and removed in PR4 because no
solver used it.

\textbf{Verification:} explicit-path tests in \texttt{test/tests/real\_space}
(Dirichlet linear profile and 2D corners, Neumann conservation and linear
profile, Robin steady state against an analytic solution, semi-infinite
and thin-slab diffusion); the R4 dual-path unit test
(\texttt{unit/src/TensorBoundaryConditionTest.C}), rewritten in PR4.

\textbf{Retired:} R5; R4 for the explicit path (implicit path: PR4).

\status{merged into \texttt{devel} (PR \#186, \texttt{fd\_boundary\_conditions},
2026-10-20).}

\pr{1a}{Torch-Native vs.\ PETSc Matrix-Free Solve Benchmark}
\textbf{Question (R6):} does the torch/PETSc hand-off erase PETSc's benefit
over a torch-native matrix-free CG? \textbf{Built:} a C++ gtest benchmark
(\texttt{unit/src/PetscTorchMatrixFreeBenchmarkTest.C}, branch
\texttt{petsc\_torch\_benchmark}, not merged) comparing the production
\texttt{MooseTensor::conjugateGradientSolve} with \texttt{KSPSolve} on a
\texttt{MatCreateShell} of the same stencil action, with zero-copy
\texttt{Vec}s; plus, at the user's request, an anisotropic case for R7 and an
assembled-\texttt{Mat} scaling sweep. The originally planned standalone
\texttt{petsc4py} script was not built.

\textbf{Results (2026-10-17).} All correctness tests pass on both CPU
(macOS arm64, this machine) and CUDA (the user's machine, confirmed via a
second run after a device-mismatch bug found in the first run was fixed
--- see below), including the zero-copy \texttt{Vec} pointer-equality
assertion on both devices.

\textit{Matrix-free timing sweep} (CG+Jacobi, \texttt{rel\_tol}=$10^{-10}$,
unpreconditioned stopping test on both legs; 32 iterations at every size
on both legs and both devices, confirming correctness is unaffected by
which leg or device is used). CPU: macOS arm64, this machine. CUDA: H200.
CUDA times are bracketed by \texttt{cudaDeviceSynchronize()}, and
\texttt{KSPSetUp} is timed separately (0.03--0.10~ms at every size). Every
CUDA leg is solved twice back to back (``first'' and ``repeat'') for the
reason given below. The CUDA table is from run~7, which was executed under
\texttt{nsys profile} (NVTX ranges on); its absolute times are therefore
slightly inflated by the profiler (launch-bound sizes $n\le512$ by up to
$\sim$2.5~ms, large sizes by $\le$1\%, compared with the unprofiled run~5
log), so the \textit{ratios} are the result, not the millisecond values.

\begin{center}
\footnotesize
\begin{tabular}{rrrr}
\toprule
$n$ & torch-native (CPU) & PETSc (CPU, zero-copy) & PETSc / torch \\
\midrule
32 & 0.53 ms & 0.36 ms & 0.68 \\
64 & 0.94 ms & 0.71 ms & 0.76 \\
128 & 3.03 ms & 2.46 ms & 0.81 \\
256 & 11.0 ms & 8.93 ms & 0.81 \\
512 & 43.4 ms & 36.2 ms & 0.84 \\
1024 & 174 ms & 147 ms & 0.85 \\
2048 & 734 ms & 616 ms & 0.84 \\
\bottomrule
\end{tabular}
\end{center}

\begin{center}
\footnotesize
\begin{tabular}{rrrrrr}
\toprule
$n$ & torch-native & torch + D2D copy & PETSc first & PETSc repeat & PETSc repeat / torch \\
\midrule
32 & 6.68 ms & 6.74 ms & 17.5 ms & 6.07 ms & 0.91 \\
64 & 6.59 ms & 6.76 ms & 6.97 ms & 6.12 ms & 0.93 \\
128 & 6.53 ms & 6.70 ms & 7.00 ms & 6.05 ms & 0.93 \\
256 & 7.43 ms & 7.63 ms & 7.05 ms & 6.12 ms & 0.82 \\
512 & 7.05 ms & 7.24 ms & 7.94 ms & 6.66 ms & 0.95 \\
1024 & 9.08 ms & 9.18 ms & 12.4 ms & 8.79 ms & 0.97 \\
2048 & 21.5 ms & 22.0 ms & 33.3 ms & 19.9 ms & 0.93 \\
4096 & 73.7 ms & 75.8 ms & 119 ms & 65.8 ms & 0.89 \\
8192 & 278 ms & 286 ms & 463 ms & 248 ms & 0.89 \\
16384 & 1094 ms & 1128 ms & 1832 ms & 974 ms & 0.89 \\
\bottomrule
\end{tabular}
\end{center}

Torch columns are the repeat solve. The first torch solve is within 0.4\% of
it for $n\ge4096$ and within 4\% for every $64\le n\le2048$. Only $n=32$
differs substantially (21.6 against 6.7~ms), because it is the process's
first CUDA work.

\begin{center}
\includegraphics[width=\textwidth]{figures/timing_sweep.pdf}
\end{center}

\textbf{CPU.} PETSc's matrix-free shell leg is faster than torch-native at
every size: 15--19\% for $n\ge128$ (32\% at $n=32$).

\textbf{CUDA, steady state.} Below $n\approx512$ every leg sits on a
$\sim$6--7~ms floor (about 200~\textmu s per CG iteration under the
profiler), so that regime measures launch latency, not the solver. For
$n\ge4096$ the three legs are bandwidth bound ($n=16384$ is
$2.7\times10^{8}$ unknowns, 58~GB of device memory in use) and the cost is a
constant per unknown and iteration: 127~ps for torch-native, 131~ps for
torch-native with the extra copy, and 113~ps for PETSc's repeat solve.
\textbf{PETSc's repeat solve is faster than torch-native at every size,
by 3--18\% (11\% at $n\ge4096$); the small-size values are within run-to-run
noise of the launch floor and were not repeated.} It is also faster than the
copy-paying torch leg (0.80--0.96$\times$, 0.87$\times$ at $n\ge4096$), which
is the fairer comparison since \texttt{cudaShellMatMult} pays that copy too.

\textbf{First solve versus repeat solve.} PETSc's \emph{first}
\texttt{KSPSolve} after \texttt{KSPSetUp} is slower than its repeat by a
one-off cost that scales as $n^{2}$: 858~ms at $n=16384$, about 3.2~ns per
unknown at every $n\ge2048$. It is the entire apparent ``PETSc is
1.68$\times$ slower'' result of run~5; in run~7 the first-solve ratio to
torch is 1.67. Torch-native has no such cost: its repeat solve equals its
first. This was established in two steps, an \texttt{nsys} trace of run~5's
sweep and a repeat solve added in runs~6 and~7 (PETSc's repeat time is
identical in the two, 974.3~ms).

\textit{Mechanism} (read from the PETSc 3.25.4 source and matched call for
call against the trace; the \texttt{jacobi.c} lines were rechecked by hand,
but nothing here has been confirmed by an experiment yet). The cost is not in
the wrapped \texttt{b} and \texttt{x} Vecs; it is in \texttt{PCJACOBI}'s
diagonal \texttt{Vec}. \texttt{KSPSetUp} does no Jacobi work because
\texttt{jac->diag} does not exist yet, which is why it takes only 0.04~ms.
\texttt{PCApply\_Jacobi} creates the diagonal lazily on the first
application, which for the unpreconditioned stopping test happens before the
CG loop starts. \texttt{MatGetDiagonal} (the shell's callback) and
\texttt{VecReciprocal} run on the GPU. Then, because \texttt{fixdiag}
(replace a zero diagonal entry by 1) defaults on and the shell is not known
to be SPD, PETSc reads the diagonal on the host (\texttt{VecGetArray}) and
loops over all $n^{2}$ entries. That forces a pinned host allocation (the
$\sim$508~ms \texttt{cudaMallocHost}), a 39~ms device-to-host copy of the
2.1~GB vector, the host loop (the $\sim$270~ms gap, an inference: no other
code runs between the two copies), and a 39~ms copy back. This replaces my
earlier guess that PETSc round-trips a mirror of the wrapped \texttt{Vec}s;
that was wrong. No CG iteration runs before the stall.

\textit{Recurrence.} From the source: not for the same \texttt{KSP} with
unchanged operators, including after wrapping new \texttt{Vec}s over the
same buffers. It does recur in full for a newly created \texttt{KSP}/\texttt{PC}. If
the operator's state changes on an existing \texttt{PC}, the diagonal is
rebuilt and the host loop and both copies return, but not the pinned
allocation. (As built in PR3, \texttt{ImplicitFDSolver} creates one
\texttt{KSP} per field, marks definite levels \texttt{MAT\_SPD}, and changes
the operator state only when the stencil changes.)

The cost belongs to \texttt{PCJACOBI}'s setup rather than to zero-copy
wrapping, so a different preconditioner would not pay it. It is avoidable
by configuration, but each option is a numerical choice:
\texttt{PCJacobiSetFixDiagonal(pc, PETSC\_FALSE)} would stop zero diagonal
entries being replaced by 1 (identical results here, where the diagonal is a
nonzero constant) and \texttt{MatSetOption(mat, MAT\_SPD, PETSC\_TRUE)}
asserts the operator is SPD. Neither has been tried. At 3.2~ns per unknown
the one-off cost is about 80\% of an entire 32-iteration torch-native solve
(4.1~ns per unknown), so creating the \texttt{KSP}/\texttt{PC} once and
reusing it across solves matters.

\textbf{Where PETSc's steady-state advantage comes from.} The matvec is the
same work in every leg: the \texttt{im2col} and \texttt{gemv} kernels total
501~ms of GPU time at $n=16384$ in each leg (46\% of torch-native's busy
time). Everything else is vector work. Torch-native spends 582~ms there in
many ATen elementwise and reduction kernels; PETSc spends about 499~ms in
cuBLAS \texttt{axpy}, \texttt{dot} and \texttt{nrm2} (first-solve trace,
which also contains the 78~ms of host copies counted separately, so this
is only a rough comparison). That fits the observed 14~ps, or $\sim$120~ms
per solve, difference, but it is a reading of kernel totals, not a
controlled experiment. Neither leg avoids synchronization: both perform
three blocking scalar reductions per CG iteration (about 3 device-to-host
copies per iteration in each leg's trace), so synchronization does not
explain the difference, and eager libtorch fuses nothing. The extra
device-to-device copy in \texttt{cudaShellMatMult} costs 3.9~ps per unknown
and iteration (3.1\%), as the copy leg shows directly.

R6 therefore resolves in PETSc's favour on both devices at steady state:
zero-copy \texttt{Vec} wrapping (Section~\ref{sec:soa}) works correctly on
CPU and CUDA, and the matrix-free PETSc leg is faster than torch-native on
both, by 15--19\% on CPU for $n\ge128$ (32\% at $n=32$) and 3--18\% on this GPU (11\% at large $n$). The condition is the
reuse requirement above; the first-solve cost is the open item.

\textit{Ill-conditioned (anisotropic, R7/PR4a) comparison} ($n=64$,
$D_x=1$, $D_y=10^{-4}$, $D_x/D_y=10^4$, same manufactured solution and
correctness check in both legs). \texttt{KSPSetUp} is called before the
timer so that PETSc's lazy \texttt{PCSetUp} (the GAMG hierarchy build) is
not charged to the first solve; ``repeat solve'' is a second solve from a
zero guess on the same \texttt{KSP} with the preconditioner reused, with the
iteration count and solution error asserted equal to the first solve.

\begin{center}
\footnotesize
\begin{tabular}{lrrrrr}
\toprule
Case & Its. & \texttt{Pmat} assembly & PC setup & First solve & Repeat solve \\
\midrule
Jacobi (CPU) & 202 & --- & 0.01 ms & 4.37 ms & 4.34 ms \\
AMG (CPU) & 134 & 0.33 ms & 3.77 ms & 18.4 ms & 18.2 ms \\
Jacobi (CUDA) & 202 & --- & 0.06 ms & 27.7 ms & 26.7 ms \\
AMG (CUDA) & 134 & 0.67 ms & 6.94 ms & 80.1 ms & 79.2 ms \\
\bottomrule
\end{tabular}
\end{center}

\begin{center}
\includegraphics[width=\textwidth]{figures/anisotropic_amg_vs_jacobi.pdf}
\end{center}

AMG needs 1.5$\times$ fewer iterations (134 against 202), identically on
both devices. That is concrete evidence for R7 that a scalar matrix-free
preconditioner cannot reach on its own. It does not translate into
wall-clock time at this size. With the preconditioner already built, an AMG
solve is 4.2$\times$ slower than a Jacobi solve on CPU (18.2 against
4.3~ms) and 3.0$\times$ slower on CUDA (79.2 against 26.7~ms). The one-time
costs (\texttt{Pmat} assembly plus PC setup) are only 4.1~ms on CPU and
7.6~ms on CUDA, so amortizing them over many solves cannot rescue AMG
here; the cost is in applying the hierarchy each iteration. Two further
points. First, $n=64$ (4096 unknowns) is deep in the CUDA latency-floor
regime: 26.7~ms over 202 iterations is again $\sim$130~\textmu s per
iteration, so the CUDA anisotropic numbers say little about AMG on a GPU.
Second, \texttt{Pmat} is a host \texttt{MATSEQAIJ}, so GAMG's hierarchy and
V-cycle run on the host and each \texttt{PCApply} presumably moves the
residual and correction across the bus. That would explain AMG being 4.3$\times$
slower on CUDA than on CPU (79.2 against 18.2~ms), but it has not been
measured. Whether AMG pays off at larger $n$ (where Jacobi-CG iteration
counts are expected to grow and AMG's to stay roughly flat) is untested.
At these costs, whether \texttt{Pmat} is built once or per step does not
change the conclusion.

\textbf{History of this comparison.} The first two CUDA runs reported
identical Jacobi/AMG iteration counts to the CPU run because the test
(\texttt{AnisotropicAMGConvergesFasterThanJacobi}) was CPU-only; a genuine
CUDA variant was added (commit \texttt{cb9368f}). Wall-clock timing was
added next (commit \texttt{24527aa}) because iteration counts alone could
not show which leg is faster. The initial wall-clock numbers charged
PETSc's lazy \texttt{PCSetUp} to the first solve; the setup/solve/repeat
split above (run~5) removes that. The CUDA matrix-free sweep was
originally reported as ``PETSc is 1.68$\times$ slower, $\sim$82~ps per unknown
and iteration unexplained''. That compared PETSc's first solve with
torch-native's first solve, and the gap was PETSc's one-off cost, not
per-iteration work. Adding repeat solves to every leg (runs~6 and~7)
withdrew that finding; the $\sim$82~ps figure and the plateau claim no
longer stand.

\textit{Assembled-\texttt{Mat} fallback scaling} (same isotropic problem as
the timing sweep, a fully assembled \texttt{Mat} for both \texttt{Amat} and
the Jacobi \texttt{Pmat}, CPU only --- not yet exercised on CUDA):

\begin{center}
\footnotesize
\begin{tabular}{rrrr}
\toprule
$n$ & assembly & solve & \texttt{Mat} storage (\texttt{nz\_used}) \\
\midrule
32 & 0.08 ms & 0.09 ms & 0 MB \\
256 & 4.32 ms & 5.44 ms & 5 MB \\
1024 & 67.8 ms & 94.6 ms & 80 MB \\
2048 & 276 ms & 394 ms & 320 MB \\
\bottomrule
\end{tabular}
\end{center}

Iteration counts match the matrix-free legs exactly at every size, so
assembly changes nothing about correctness or convergence, only cost. At
$n=2048$ the assembled matrix occupies 320~MB on its own. Once assembled,
its solve (394~ms) is faster than either matrix-free CPU leg (616 and
734~ms); with assembly included (670~ms) it falls between them. For this
constant-coefficient 5-point stencil on CPU, the evidence for the
matrix-free default (Section~\ref{sec:decisions}, item 5) is therefore
storage and device residency, not speed. \textbf{Peak RSS is deliberately
not tabulated:} \texttt{ru\_maxrss} is a process-wide high-water mark and
this sweep runs after the matrix-free sweep in the same process (which had
already reached 1121~MB), so the log's RSS column does not measure the
assembled path's own footprint. An earlier revision of this section
compared those numbers; that comparison was not valid and is withdrawn.

\textbf{Fix found and applied mid-verification:} the first CUDA run
(before the fixes above) failed both new CUDA-only tests with a device
mismatch (\texttt{conv2d}'s weight tensor was unconditionally CPU-resident
in \texttt{PeriodicStencilOperator}, so applying it to a CUDA input threw).
Fixed by giving the stencil operator classes an optional \texttt{device}
parameter (commit \texttt{19341d7}); a second CUDA run with the fix passed
all 8 tests. This is recorded here, rather than only in commit history,
because it is exactly the kind of PETSc/torch hand-off correctness issue
Section~\ref{sec:soa}'s zero-copy requirement is meant to get right, and it
only surfaced once real CUDA hardware exercised the code.

The full raw logs from both machines are reproduced below for direct
inspection, rather than only the extracted tables above.

\lstinputlisting[caption={Full CPU run log (\texttt{petsc\_torch\_benchmark\_cpu.log}), run~5 (setup/solve/repeat timing), 6/6 tests passing.}]{logs/petsc_torch_benchmark_cpu.log}

\lstinputlisting[caption={Full CUDA run log (\texttt{petsc\_torch\_benchmark\_cuda.log}), run~5 on an H200, all 9/9 tests passing. Its \texttt{DISABLED\_CudaTimingSweep} lines time first solves only and are superseded by the repeat-solve log below; the other tests' lines (including the anisotropic CUDA test) are current.}]{logs/petsc_torch_benchmark_cuda.log}

\lstinputlisting[caption={CUDA timing sweep with first and repeat solves for every leg (\texttt{petsc\_torch\_benchmark\_cuda\_sweep.log}), run~7 on an H200 under \texttt{nsys profile}, 1/1 tests passing.}]{logs/petsc_torch_benchmark_cuda_sweep.log}

\status{implemented and verified on both CPU and CUDA, all 9 tests,
2026-10-17 (commits \texttt{12016ae}, \texttt{41c2e31}, \texttt{19341d7},
\texttt{cb9368f}, \texttt{24527aa}, \texttt{c77191a}, \texttt{5c38c75},
\texttt{d0f8d60}, \texttt{3ac2c7c} on branch \texttt{petsc\_torch\_benchmark}).
R6 and R7 are both evidenced on CPU and CUDA. At steady state PETSc's
matrix-free leg is faster than torch-native on both devices (CUDA: 11\% at
large $n$). Its first \texttt{KSPSolve} carries a one-off, $n^{2}$-scaling
cost (3.2~ns per unknown) traced from the PETSc source to
\texttt{PCJACOBI}'s lazily built, host-checked diagonal (not yet confirmed by
an experiment). The independent-reviewer-agent pass was run together with
PR1b's on 2026-10-18 and its findings are fixed. Not merged (spike, tag
\texttt{pr1b-spikes}); its role passed to PR3, which creates one
\texttt{KSP} per field and reuses it.}

\pr{1b}{Multigrid Spikes and the Hybrid-Design Decision}
\label{sec:pr1b}
These are spikes, not maintained solvers, on branch \texttt{petsc\_torch\_benchmark}.
Unless stated otherwise every run is 2D, periodic, with the SPD operator
$-\Delta u + \sigma u$ ($\sigma = 1$). The CUDA runs were done by the user on an
H200 with PETSc 3.25.4 and libtorch 2.10. The table lists the tests in the
order they were run.

\begin{longtable}{>{\bfseries}p{0.03\textwidth} p{0.25\textwidth} p{0.4\textwidth} p{0.23\textwidth}}
\toprule
\# & Test (commit) & Result & Led to \\
\midrule
\endhead
1 & PR1a Jacobi-CG benchmark, CPU + CUDA (\texttt{12016ae}\ldots\texttt{3ac2c7c}) & Zero-copy holds on both devices. At steady state PETSc is 15--19\% faster on CPU ($n\ge128$) and 3--18\% faster on CUDA. The first CUDA solve is slow & Why? \\
2 & First-solve source trace (\texttt{27a8807}) & \texttt{PCJACOBI} lazily builds its diagonal and checks it on the host; the \texttt{Vec} wrapping is not the cause & Jacobi-CG is $O(n)$ iterations: is torch GMG feasible? \\
3 & Torch GMG, emulated slabs (\texttt{c3364ec}) & RBGS V-cycle MG-PCG takes 11 iterations at $n=64..512$, identical bit for bit on 1/2/4 slabs & Real MPI? \\
4 & Torch GMG, real MPI (\texttt{34e10ff}) & 10 iterations at np $=1/2/4/8$ with the same error. Dropping one halo update breaks it & DMDA instead? \\
5 & DMDA decomposition spike (\texttt{18d5886}, unmerged) & A DMDA forced onto Marlin's partition matches ownership and ghosts and aliases the data zero-copy on 1--8 ranks & R1 feasible \\
6 & Torch V-cycle CUDA sweep (\texttt{e566e3c}) & $\sim$2~ms eager launch floor. MG-PCG takes 11 iterations at every $n$ and beats Jacobi-CG by 4.5$\times$ to $>$360$\times$ & Does PETSc GMG match? \\
7 & PETSc KSPCG+PCMG on torch shells, CPU (\texttt{610802e}) & 9 iterations, grid independent & Time it on the GPU \\
8 & PCMG CUDA sweep (\texttt{6698261}) & 8--9 iterations against torch's 11, so the repeat solve is 0.70--0.92$\times$ torch's. That compares different smoothers: per iteration and per V-cycle PETSc is 13--38\% slower at $n\ge2048$ (44 vs.\ 32~ms per V-cycle at $n=8192$), with the same $\sim$2~ms floor. \texttt{KSPSetUp} costs 1.5--1.8$\times$ a solve at $n\ge2048$ and 0.6--1.0$\times$ at $n=512$--1024 & Paid every Picard step? \\
9 & Picard update-policy sweep (\texttt{bd00ff4}, \texttt{665ac06}, \texttt{be91418}) & A state bump triggers only a \emph{partial} lazy rebuild (finest smoother and Jacobi diagonals; review finding), and the coarse LU was refreshed every step under every policy. Stale setups can go indefinite, so retry is mandatory. Threshold reuse saves 39\% at $\beta=3$. At $\beta=9$ every policy fails the DTOL test ($\|b\|$ reference with a warm start) & Stopped tuning the test; audited our PETSc usage \\
10 & \texttt{-log\_view} audit & Every host$\leftrightarrow$device copy is in \texttt{PCSetUp}: the Jacobi zero-diagonal check, Chebyshev's GMRES estimate, and Chebyshev's random right-hand side, which is built on the host (\texttt{KSPSetNoisy\_Private}) regardless of \texttt{MAT\_SPD}. \texttt{PCApply} and \texttt{MatMult} make no copies. The log on disk is an excerpt, so this is the recorded reading rather than a re-checkable one & Tell PETSc the operator is SPD \\
11 & \texttt{MAT\_SPD} variant (\texttt{2ded6f0}) & Setup 724$\to$408~ms at $n=8192$; nothing else changes (the $n=256$ pair is confounded by run order). The Chebyshev estimate remains, including its host-built random right-hand side & Gershgorin bounds deferred; hybrid design adopted \\
\bottomrule
\end{longtable}

Logs: \texttt{doc/roadmap/logs/torch\_multigrid\_cuda\_vcycle\_sweep.log},
\texttt{petsc\_multigrid\_cuda\_sweep.log} and
\texttt{petsc\_multigrid\_cuda\_sweep\_log\_view.log} (gitignored).

\textbf{Not merged as-is (decided 2026-10-18).} Apart from PR1a's
torch-native leg, these tests exercise spike code that lives in the test
files, so as regression tests they would mostly check PETSc and torch. The
branch is preserved at tag \texttt{pr1b-spikes}. Their assertions carry
forward into PR3, rewritten against the production solver: zero-copy
pointer equality (CPU, CUDA); grid-independent iteration counts;
\texttt{MAT\_SPD} leaving the solution and iteration count unchanged;
rank-count-independent iteration count with a dropped-halo negative check
(PR2b); and every preconditioner-update policy reaching the same iterate,
with rebuild-and-retry on a negative converged reason.

\textbf{Independent review (2026-10-18).} No defects found in the stencils,
transfers, red-black colouring, halo exchange, PETSc array get/restore or
timer synchronisation. The review found the partial-rebuild, coarse-LU and
\texttt{MAT\_SPD}-reset mechanisms and the number corrections now reflected
above and in decision~3. The test-code findings were fixed in
\texttt{016081d}: the MPI test now asserts the iteration count at every
rank count, a leaked shell Mat, and stale comments. The CPU tests and MPI
np~$=1/2/4/8$ pass after the fixes.

\status{implemented and independently reviewed; findings addressed. Spikes
only, not merged; not on the production path.}

\pr{2}{DMDA/Decomposition Compatibility Spike}
\textbf{Question (R1):} can a DMDA be forced onto Marlin's decomposition,
with owned ranges matching the torch tensors exactly? \textbf{Answer:} yes.
The spike (\texttt{18d5886}, unmerged) passes on 1/2/3/4/6/8 ranks using
only public \texttt{DomainAction} getters; no \texttt{DomainAction} addition
was needed. In production, \texttt{PetscFD::gatherPartition} and
\texttt{buildHierarchy} do this and check ownership on every run, periodic
and non-periodic (\texttt{DM\_BOUNDARY\_NONE}) axes alike, tested by
\texttt{PetscFDHierarchyTest.FineDMDAMatchesMarlinBlock}.

\status{spike done, not merged; superseded by \texttt{PetscFDHierarchy} (PR3, PR4). R1 retired.}

\pr{2b}{Multi-Rank PCMG on Torch Shell Levels (Spike)}
\textbf{Scope:} Run PR1b's PCMG problem on real MPI ranks through
\texttt{DomainAction}. Settle the level halo exchange (per-level DMDA
vs.\ Marlin's own exchange), the coarse gather
(\texttt{PCREDUNDANT}/\texttt{PCTELESCOPE}), and how to coarsen thin slabs.
Add a one-rank 3D check.

\textbf{Verification:} Iteration count and solution must be independent of
the rank count (np $=1/2/4/8$). Dropping one halo update must break
convergence.

\textbf{Retires:} R8 (2D periodic; non-periodic axes still open). \status{spike done
(\texttt{46ecff3}, \texttt{cd56180}, \texttt{6f9b8ee}, \texttt{99bc4ec}, \texttt{b3dcd72}). Own MPI halo
exchange in each shell callback, assembled MPIAIJ coarsest level solved by
\texttt{PCREDUNDANT}, and halving while every owned extent stays even.
On CPU at np $=1/2/4/8$: 9 iterations at $n=64/128/256$, and at $n=96$
(coarse $n=12$), at every rank count, with relative error $\le 2.1\times10^{-13}$.
Caveat: at $n=112$ on $2\times4$ ranks a slab turns odd, so the hierarchy
stops early (coarse $n=28$, 7 iterations): the redundant coarse problem
grows with the rank count, which needs agglomeration or a rule on grid sizes.
Dropping the fine halo gives relative error 0.61--0.76 while KSP still reports
convergence, so PR3's tests need an error check, not only a convergence check.
\begin{itemize}[leftmargin=*]
\item \textbf{\texttt{PCTELESCOPE}.} The default telescope mode extracts rows
with \texttt{MatCreateSubMatrices}, so it needs an assembled \texttt{Mat} and
only wraps the assembled coarse level, not the torch shells. Gathering to one
rank with an inner preonly + LU gives the same iteration count as
\texttt{PCREDUNDANT} at every rank count (it solves the same coarse problem,
on one rank instead of all). With an inner \texttt{PCGAMG} (one AMG V-cycle
instead of an exact coarse solve; a preconditioner choice) the counts are
9--10, and 11 on the shallow $n=112$, $2\times4$ hierarchy. Telescope removes
the redundant work on the other ranks but not the coarse-size growth; that
still needs a rule on grid sizes or an inner solver that restores depth.
\item \textbf{Per-level DMDA halo.} A DMDA per level, forced onto Marlin's
partition (explicit \texttt{lx}/\texttt{ly}, axes swapped so DMDA's x-fastest
order is torch's row-major $[m_x, m_y]$, ranks renumbered to match; ownership
checked against Marlin). Ghosts are filled by \texttt{DMGlobalToLocal} on Vecs
placed onto the tensors' storage. Solutions and iteration counts are
bitwise identical to our own exchange at np $=1/2/4/8$. For the halo alone the
two are equivalent; neither fills non-periodic boundary ghosts
(\texttt{DM\_BOUNDARY\_GHOSTED} leaves them to Marlin's
\texttt{TensorBoundaryCondition}, as \texttt{exchangeGhostLayers} does). The DMDA's
value is the hierarchy below. Timings were not measured.
\item \textbf{PETSc-owned hierarchy} (\texttt{b3dcd72}). Only the fine level is a
torch shell. The coarse DMDAs and their partitions come from \texttt{DMCoarsen},
the coarse levels are rediscretised and assembled with
\texttt{MatSetValuesStencil} (scaled by $4^k$ because PCMG restricts with
$P^T$; a preconditioner-level convention), and the transfers come from the DM or
are assembled from our stencil. The depth no longer depends on the rank grid:
$n=112$ on $2\times4$ keeps 4 levels (coarse $n=14$), 9 iterations. Iterations at
np $=1/2/4/8$, $n=64/96/112/256$:
cell-centred bilinear (our stencil, assembled) 9 at every case, equal to the
single-rank torch-shell hierarchy; \texttt{DMDA\_Q1} 9--10; \texttt{DMDA\_Q0}
17--19. The fully automatic path (\texttt{KSPSetDM} + \texttt{KSPSetComputeOperators},
PCMG coarsens and builds Q1 itself) gives the same counts as the hand-built Q1
hierarchy. PETSc's coarse partition is chosen for vertex-based transfers, so our
cell-centred interpolation is assembled with global column indices (a needed
coarse cell can lie outside the one-cell ghost frame). \textbf{Decided (user,
2026-10-18):} adopt this hierarchy (Section~\ref{sec:hybrid}) with the
cell-centred transfer (exact for our cell-centred grid; Q1 is node-centred).
\item \textbf{3D, one rank.} 7-point operator, the same transfers applied axis
by axis, assembled SeqAIJ coarse level: 8 iterations at $n=16/32/64$ with both
\texttt{PCREDUNDANT} and telescope + LU, relative error $\le 8.7\times10^{-12}$.
\end{itemize}
Not covered: multi-rank 3D, non-periodic axes, GPU, timings.}

\pr{2c}{Non-SPD Operator with a Torch \texttt{PCSHELL} Smoother + Fusion Measurement (Spike)}
\textbf{Scope:} Advection--diffusion solved with KSPGMRES/FGMRES + PCMG.
Compare a torch symmetric red-black Gauss-Seidel \texttt{PCSHELL} smoother
against PETSc's own smoother. Time the same smoother eager, fused through
\texttt{TracedComputeSequence}, and CUDA-graph captured.

\textbf{Verification:} Iterations and solution against a reference solve;
timings reported, not assumed.

\textbf{Retires:} R9 (partly), R10 (CPU only so far). \status{CPU spike done
(\texttt{40f6977}). Periodic $-D\nabla^2u + \mathbf{v}\cdot\nabla u + \sigma u$,
$\mathbf{v}$ at $30^\circ$, upwind and central advection, fine-grid cell P\'eclet
0.1/1/10; KSPGMRES + PCMG on the PETSc-owned hierarchy; $n=64/128/256$.
\begin{itemize}[leftmargin=*]
\item Chebyshev/Jacobi works at Pe $\le 1$ upwind (10--21 its) and fails at upwind
Pe $=10$ and central Pe $\ge 1$ (500 its).
\item Gauss--Seidel smoothers are robust for upwind at every Pe: torch symmetric
red-black GS (\texttt{PCSHELL}, fine level) 10--29 its, PETSc symmetric
\texttt{PCSOR} 10--25 its.
\item Central differencing at Pe $\ge 1$ is solved by no smoother here (Pe $=1$
breaks down for $n\ge128$; Pe $=10$ gives NaN). Expected cause, not diagnosed:
the coarse levels see Pe doubled per level and lose the M-matrix property. A
discretisation/coarse-operator question, not a smoother one.
\item The traced smoother equals eager to the bit; on CPU it gives no speed-up
(0.575 vs 0.571~ms per PCApply at $n=256$; 9.4 vs 9.0~ms at $n=1024$).
\end{itemize}
\textbf{Central advection at high Pe is deferred (user, 2026-10-18)} until a
flow-coupled problem needs implicit advection: no current target physics has an
advection term; high cell Pe would come from solute in a flowing melt
($D\sim\SI{e-9}{\metre\squared\per\second}$, Pe $\sim1$--100 at \si{\micro\metre} grids). Options
then, each a separate numerical-method decision: explicit advection in $b$
(implicit matrix stays SPD; CFL limit); an upwind-type implicit scheme; or a
split (defect correction): enough of the advection implicit to keep the operator
well conditioned for multigrid, the remainder explicit in $b$, Picard-iterated.
At Picard convergence this is the fully implicit solution of the target scheme
(a truncated loop is a lagged scheme), and it converges only while the explicit
remainder is small relative to the implicit part.
\textbf{CUDA sweep (H200, \texttt{DISABLED\_CudaSmootherFusionSweep}), ms per
\texttt{PCApply} of the fine RBGS PCSHELL:}
\begin{center}
\begin{tabular}{rrrrr}
\toprule
$n$ & eager & traced (FuseGraph) & CUDA graph (eager) & CUDA graph (traced) \\
\midrule
256 & 0.185 & 0.060 & 0.086 & 0.033 \\
512 & 0.177 & 0.059 & 0.129 & 0.046 \\
1024 & 0.346 & 0.143 & 0.284 & 0.131 \\
2048 & 1.261 & 0.518 & 1.193 & 0.510 \\
4096 & 4.489 & 1.859 & 4.465 & 1.901 \\
8192 & 17.301 & 7.260 & 17.464 & 7.493 \\
\bottomrule
\end{tabular}
\end{center}
(Run at \texttt{65109ae}; the earlier three-mode runs agree within 2\%.)
Tracing gives 2.4--3.2$\times$ at every size (fused kernels, less memory
traffic); capturing eager in a CUDA graph only removes launch overhead
(2.1$\times$ at $n=256$, none for $n\ge2048$). Capturing the traced smoother
combines both: 1.8$\times$ over traced at $n=256$ (0.033~ms, 5.6$\times$ eager),
1.3$\times$ at $n=512$, equal at $n=1024$--2048, and 2--3\% slower at $n\ge4096$
(the copy into the captured input buffer). This answers R10: fusion is worth
having on CUDA. \textbf{Graph capture is not adopted (user, 2026-10-18):} it
gives nothing on large systems, and on small ones the run times are already low
enough that the gain rarely matters.
All four modes match the CPU eager result to $\le5.2\times10^{-16}$ relative
(the first run's check failed only because its input came from the CUDA random
generator; fixed in \texttt{63f0df1}).}

\pr{3}{\texttt{ImplicitFDSolver}: Periodic, Matrix-Free, No Picard}
\textbf{Built} (design name \texttt{PetscFDSolverBase}):
\texttt{ImplicitFDSolver : TensorSolver, IterativeTensorSolverInterface}
(\texttt{include/tensor\_solver/ImplicitFDSolver.h},
\texttt{src/tensor\_solver/ImplicitFDSolver.C}) and the \texttt{PetscFD}
utilities (\texttt{src/utils/PetscFDHierarchy.C}); the solver block takes
\texttt{buffer}, per-field \texttt{operators} and \texttt{source},
\texttt{linear\_solver}, \texttt{l\_tol}, \texttt{l\_max\_its} and MOOSE's
\texttt{petsc\_options*} (prefix \texttt{<field>\_}).
\texttt{TimeDerivativeOperator} (planned for PR5a) was pulled forward. The
operators must be \texttt{LinearFDOperator}s acting on the solved field and
SPD or SPSD; a singular field's right-hand side must have zero mean up to
round-off. Section~\ref{sec:architecture} describes the backends.
Not built from the PR3 design: the Chebyshev analytic-bounds hook; a
pointer-equality zero-copy test in the production tests (zero-copy is by
\texttt{VecPlaceArray} on the solver's tensors and was checked by
\texttt{-log\_view} on H200). The DTOL default with warm starts was left to PR3b
(decision~9).

\textbf{Verification:} \texttt{test/tests/implicit\_fd}: Poisson MMS
(singular, zero-mean), backward Euler against the exact discrete
amplification of a Fourier mode (2D/3D, 3/5-point, both backends, CUDA),
grid-independent PCMG iterations, uneven rank ownership on 3 ranks, PETSc
options reaching the KSP, failed-solve cutback, and setup errors;
\texttt{unit/src/PetscFDHierarchyTest.C}.

\status{merged into \texttt{devel} (PR \#188, \texttt{implicit\_fd\_solver}, issue \#184,
2026-10-20), independently reviewed, findings addressed. Tests pass on CPU at 1, 2 and 4 ranks,
MPS and CUDA (H200). Physics record (CPU, $l_{tol}=10^{-12}$, $n=32\to256$): Poisson MMS
observed order 2.00 (3-point) and 4.00 (5-point) for both backends, identical errors;
backward Euler against the exact spectral BE solution, order 2.00 / 4.00, while the FFT
integrator (AdamsBashforthMoulton AB1, the FFTSemiImplicit update) matches that solution to
$10^{-15}$. H200, $n=256$: 7 PCMG iterations, no host/device copies in \texttt{KSPSolve}
with Jacobi smoothing on every level.}

\pr{3b}{Implicit Biharmonic, PFHub 1a, Smoother and Coarse-Matrix Fixes}
\textbf{Built} (commits tagged ``(PR3b)'' on \texttt{implicit\_ch\_biharmonic}):
\begin{itemize}[leftmargin=*]
\item \texttt{FiniteDifferenceBiharmonic}
(\texttt{src/tensor\_computes/FiniteDifferenceBiharmonic.C}), the composed
3- or 5-point Laplacian squared, as a \texttt{LinearFDOperator}; periodic
only in \texttt{ImplicitFDSolver} (an error with boundary faces).
\item The implicit FD PFHub 1a input
(\texttt{benchmarks/01\_spinodal\_decomposition/1a\_implicit\_fd.i}):
backward Euler with $\partial_t c$ and the biharmonic implicit,
$M\nabla^2 f'(c^n)$ explicit as the source; XDMF output, solve timing, and a
plot in the benchmark notebook. \texttt{TensorIntegralChangePostprocessor}
works with ghost layers on several ranks.
\item The \texttt{smoother} parameter: PCMG's Chebyshev by default, or damped
Jacobi with the Gershgorin weight; coarse Chebyshev bounds re-estimated when the
operators change; the coarsest level solved by LU (decision~8).
\item No PETSc divergence test (decision~9).
\item Review fixes for \#189: the biharmonic kernel is built once; the Jacobi
diagonal asserts a positive centre; a 2-rank
\texttt{integral\_change} test; the notebook skips a missing implicit FD output.
\item Coarse matrices preallocated from the stencil's columns (the DMDA box
preallocation stored 729 entries per row for 73 used with the 3D 5-point
biharmonic); stencil entries that wrap onto the same coarse column are
summed; \texttt{-<field>\_pc\_mg\_levels} is an error.
\item Size-scaling script \texttt{benchmarks/03\_implicit\_fd\_scaling}
(CUDA, MPS or CPU; $512^2$ to $8192^2$; timing only, not PFHub results).
\end{itemize}

\textbf{Verification:} implicit biharmonic Fourier mode against the exact
discrete amplification in 2D and 3D, 3/5-point, both backends, CUDA and MPS
(MPS at $l_{tol}=10^{-5}$); smoother tests (default type, Chebyshev after a
$dt$ jump, damped Jacobi weight); large-$dt$ solve without a divergence stop;
coarse-assembly unit
tests (\texttt{AssembledOperatorAddsWrappedOffsets},
\texttt{StencilMatrixStoresOnlyTheStencil}).

\status{merged into \texttt{devel} (PR \#189, \texttt{implicit\_ch\_biharmonic},
2026-10-26). PR4 is built on it. CUDA run of the new Chebyshev default: not recorded here. Scaling results:
not recorded here.}

\pr{4}{Boundary Conditions in \texttt{ImplicitFDSolver}}
\textbf{Built:} \texttt{ImplicitFDSolver} solves on non-periodic domains with
\texttt{TensorRobinBC} faces on both backends. No boundary method was added
to \texttt{TensorBoundaryCondition}: the fine level reuses the ghost fill
through the affine split, and the matrix entries come from
\texttt{TensorRobinBC::ghostCoefficientOnU0}
(Section~\ref{sec:architecture}, Boundary-Condition Dual-Use Invariant).
The planned \texttt{applyGhostContribution()} was dropped, and the unused
\texttt{stencilContribution()} of PR1 was removed.

\textbf{As built} (decisions by the user, 2026-10-22):
\begin{itemize}[leftmargin=*]
\item The Robin value $c$ (Dirichlet \texttt{value}, Neumann \texttt{flux}) varies along the
face, from a MOOSE \texttt{Function} at the face position and the solve time, or from a
buffer whose first owned layer holds the face value (tangential halo exchanged). $a$ and
$b$ stay scalar. The implicit solve evaluates functions at $t^{n+1}$; a buffer holds what its compute wrote at the start of the time step, $t^n$.
\item Fine level: the operator action is affine; each solve takes $b \leftarrow b - A(0)$
and applies $A(x) - A(0)$, so the discretisation is unchanged. The torch PCG Jacobi
diagonal includes the ghost coefficient in the cells next to a face.
\item Coarse levels (preconditioner only): ghost layer $l$ is folded onto mirror cell $l$
with the homogeneous factor $-(a/2 - b/s)/(a/2 + b/s)$, $s = (2l+1)h$; the cell-centred
interpolation folds the same way and restriction stays $P^T$. \texttt{DMDA\_Q0}
interpolation keeps \texttt{DMCoarsen} cell-centred on non-periodic axes. This answers the
deferred question on coarse BC rows: rediscretised.
\item Scope: time derivative and Laplacian; 3-point with Dirichlet, Neumann and Robin,
5-point with Dirichlet and Neumann (5-point Robin is nonsymmetric and an error). Robin
with $a/(a/2 + b/\Delta x) < 0$ is an error. Biharmonic with no-flux, and PFHub 1b, are a
follow-on PR (PR5).
\end{itemize}

\textbf{Verification:} \texttt{test/tests/implicit\_fd/bc\_mms.i} (both
backends, 3/5-point, steady and backward Euler),
\texttt{grid\_independence\_bcs}, setup-error tests;
\texttt{PetscFDHierarchyTest} (folded coarse rows and interpolation, Robin
face, Gershgorin bound with faces); \texttt{TensorBoundaryConditionTest}
(fill dependence equals the solver's ghost coefficients).

\textbf{Retired:} R4 for the implicit path.

\status{implemented on branch \texttt{implicit\_fd\_bcs}; reviewed by a pre-push review agent,
findings addressed; tests pass on CPU serial and multi-rank (user); rebased onto the
\#189 head \texttt{0c38df2} (Chebyshev default smoother), CPU serial tests pass; PR \#191 open. MMS on the unit square with Dirichlet, Robin and two Neumann faces (function and
buffer values) matches an independent sparse direct solve of the same discretisation to
$10^{-6}$ relative for both backends; observed order 1.97--1.99. PCMG iterations
($l_{tol}=10^{-10}$, CPU) at $n = 64, 128, 256$: 7, 6, 6 with boundary faces, against 9, 9, 9
periodic, the same with either smoother. Crank's semi-infinite and thin-slab diffusion
solutions (the explicit \texttt{real\_space} benchmarks, with the implicit solver): error
$5.8\times10^{-3}$, $1.7\times10^{-3}$, $6.7\times10^{-4}$ at 25, 100, 400 substeps, first order
down to the semi-discrete floor $3.2\times10^{-4}$, both backends, golds from an independent
dense solve. MPS and CUDA runs of the PR4 tests: not recorded here.}

\pr{4a (contingent)}{AMG-Preconditioned Mixed Mode}
\textbf{Scope:} an assembled \texttt{Pmat} from the same
\texttt{LinearFDOperator} stencils (\texttt{stencilAt()}), passed as
\texttt{KSPSetOperators(ksp, Amat, Pmat)} with the matrix-free shell kept as
\texttt{Amat}, so that \texttt{-<field>\_pc\_type hypre}/\texttt{gamg} can
be used; reassembled when the stencil changes. A fully assembled
\texttt{Amat} is a last resort beyond that.

\textbf{Rationale:} R7. PCMG with point Jacobi may converge slowly on
strongly anisotropic or heterogeneous coefficients; AMG coarsens on the
matrix entries. PR1a measured AMG at 1.5$\times$ fewer iterations at
$D_x/D_y = 10^4$, $n=64$, but slower in wall time.

\textbf{Verification:} iterations and wall time against PCMG on an
anisotropic case, with \texttt{Pmat} assembly cost.

\status{proposed, contingent: build only if PR5/PR6 coefficients show PCMG
is insufficient.}

\pr{5}{Variable Coefficients and the No-Flux Biharmonic}
\textbf{Scope} (user, 2026-10-25: the former PR4b bundled with variable
coefficients as one capability upgrade):
\begin{itemize}[leftmargin=*]
\item \texttt{FiniteDifferenceDivGrad} (user-chosen name), a real-space
$\nabla\cdot(D\nabla u)$ (times \texttt{factor}) implementing \texttt{LinearFDOperator}, with $D$
from a buffer; usable for diffusion, Poisson with $\varepsilon(x)$ and variable
mobility. Explicit and implicit share one face-flux formula.
\item \texttt{FiniteDifferenceBiharmonic} with non-periodic faces in
\texttt{ImplicitFDSolver} (no-flux Cahn-Hilliard); the PFHub 1b input, and 2b
on the \texttt{pfhub\_2a\_implicit\_fd} branch. 6a stays in PR5a (it needs a
steady field coupled to a transient one).
\item The torch-PCG Jacobi diagonal from the folded rows, so box stencils get
their corner terms.
\end{itemize}
The lagged nonlinear remainder needs no new path: any explicit compute
writing the \texttt{source} buffer already enters $b$ (PR3b's 1a input does
this).

\textbf{Design:} a per-row stencil on \texttt{LinearFDOperator} for
variable coefficients, used by the folded assembly and diagonal, with the
constant-stencil path unchanged; $D$ restricted to the coarse levels. Folding
the composed 13-point biharmonic stencil with the Dirichlet or Neumann ghost
factor ($\mp1$) gives exactly $L_{bc}L_{bc}$ (1D Neumann row 0:
$[2,-3,1]$; Dirichlet rows 0, 1: $[10,-5,1]$, $[-5,6,-4,1]$), so the coarse
levels need no new folding; the fine level applies the two Laplacians in turn.

Decided (Section~\ref{sec:decisions}): the face coefficient, the sign of $D$,
the biharmonic faces, $D$ at a domain face and the coarse $D$.

\textbf{Verification:} MMS with no-flux faces and with spatially varying
$D$ against independent SciPy solves; mass conservation; a unit test of the
folded biharmonic against \texttt{MatMatMult}$(L_{bc}, L_{bc})$ and of the
folded diagonal against \texttt{MatGetDiagonal}; PCMG iterations against
$n$. A parallel-plate capacitor (user, 2026-10-25): two dielectric slabs
($\varepsilon_r$ 1 and 100) between plates at a potential difference $V$,
zero-flux top and bottom so the problem is 1D; the plate charge against the
series capacitance $C/A = \varepsilon_0 / (d_1/\varepsilon_{r1} +
d_2/\varepsilon_{r2})$, a uniform displacement and the potential split
between the slabs. Pass/fail with Dirichlet faces as the plates; metal plates
as $\varepsilon_r = 10^9$ inside the domain are run and their error and
iterations recorded (trouble there is evidence for PR4a, not a PR5 blocker).

\textbf{Capacitor results} (2026-10-26, \texttt{test/tests/implicit\_fd/capacitor.i},
$L = 1$\,mm, $V = 10$\,V, $d_1 = L/4$, CPU, \texttt{l\_tol} $= 10^{-13}$):
\begin{itemize}[leftmargin=*]
\item \emph{Dirichlet plates} (gated, $n = 32$): the charge at both plates and inside
slab 1 and the potential match the exact values to $10^{-14}$ with torch PCG and PCMG,
on 1 and 4 ranks; the arithmetic face mean is $6\times10^{-2}$ off. PCMG takes 21, 23,
23, 22 iterations at $n = 32$ to 256; torch Jacobi PCG 140, 278, 534 at $n = 32$ to 128.
\item \emph{Metal plates} ($\varepsilon_r = 10^9$, thickness $L/8$, recorded): the charge
inside slab 1 is $2\times10^{-9}$ (torch) and $5\times10^{-10}$ (PCMG) off at $n = 32$;
the potential $7\times10^{-9}$ and $2\times10^{-9}$ (relative to $V$). The plate-face
charges are $10^{-5}$ to $10^{-4}$ off: the potential drop across a plate is
$6\times10^{-10}$\,V, so round-off in $\phi$ dominates the face flux there (a
NumPy direct solve of the same 1D system shows the same, $10^{-6}$ to $10^{-5}$,
while its mid-slab flux is exact to $10^{-15}$). Torch Jacobi PCG takes 94, 188, 379
iterations at $n = 32$, 64, 128; PCMG takes 1763, 4597, 9565 (2, 5, 13\,s), growing
with $n$ and slower than Jacobi. The Jacobi-scaled condition number of the 1D system is
385, 1590, 6460, so the diagonal scaling removes the $10^9$ contrast; the coarse levels
of PCMG, with $D$ averaged across the metal--dielectric faces, do not. Evidence for
PR4a (or a Galerkin coarse operator) when such contrasts matter.
\end{itemize}

\textbf{DivGrad face types and a changing $D$} (2026-10-26,
\texttt{divgrad\_mms.i}): a Robin face (steady and backward Euler) and a coefficient
switched to a jump of 100 at the start of the second step match SciPy sparse solves of
the same discretisation to $10^{-12}$ relative, with torch PCG and PCMG on 1 and 4 ranks.
The explicit computes see the time at the start of the step (decision~11), so the switch
applies from the second solve; PCMG takes 11, 19, 18 iterations at $n = 32$. With the
coarse levels not reassembled after the switch the second solve fails.

\status{implemented and verified on branch \texttt{implicit\_fd\_pr5} (on top of PR4),
tag \texttt{pr5-variable-coefficient}, 2026-10-26. The milestone review (numerics,
parallel, tests, docs, verifier) and the pre-push review found: a hang when $D$ changed on
some ranks only (reassembly is now collective), a PETSc segfault on a zero diagonal (now an
error), the biharmonic face pairs holding only for face data linear along the face (fixed
with the approved two-stage operator, for every use), and a per-iteration coefficient check
(now once per change); all fixed with tests. CUDA runs in PR CI. The full PFHub 1b run and
the 2b branch rebase are the user's.}

\pr{5a}{Steady Field Coupled to a Transient Field}
\textbf{Scope:} \texttt{TimeDerivativeOperator} is built (PR3) and steady
fields solve (Poisson tests). Remaining: a steady field (e.g.\ a Poisson
potential) and a transient field in one \texttt{ImplicitFDSolver}, coupled
through lagged sources; re-solving the steady field against current
sources is left to PR6.

\textbf{Verification:} a coupled two-field case against an independent
solution.

\textbf{Decisions (user, 2026-10-26).}
\begin{itemize}
\item \textbf{D8, coupling order.} Every source comes from $t^n$, whatever the
  order of \texttt{buffer}. A Gauss-Seidel sweep in \texttt{buffer} order
  (results depend on the order) and solving the transient fields before the
  steady ones were proposed and declined; a block-coupled solve
  (PCFIELDSPLIT with cross-field operators) is deferred until after Picard
  (PR6).
\item \textbf{D9, the mean of the solution.} When the operator maps the
  constants to $\lambda > 0$ times the constants (A symmetric), the summed
  equations fix $\bar u = \bar b / \lambda$ exactly; the solver detects this
  from $A\mathbf{1}$ when the operator is applied and sets the mean of the
  Krylov solution to $\bar b/\lambda$ after each solve, for every field
  (nothing to mark). Found by the coupled test: PCMG's SOR smoothing (and Jacobi
  next to boundary faces) leaves the mean wrong by about $l_{tol}\|u\|$ per
  step, which made the right-hand side of the singular steady field
  incompatible at $l_{tol} = 10^{-8}$. Widening the singular-field mean check
  to $l_{tol}$ was proposed and withdrawn (it hides the drift and still fails
  as the drift accumulates). MOOSE has no automatic equivalent: PETSc's
  transpose null space silently projects the right-hand side, and the
  navier\_stokes \texttt{NSPressurePin} (\texttt{average}) is an opt-in
  offset to a user target.
\item \textbf{D10, the singular-field mean check.} The round-off allowance for
  the mean of a singular field's right-hand side is scaled by the largest
  $\sum|b|$ of its solves so far, not the current one. Found by the milestone
  review: \texttt{steady\_transient.i} run to 200 steps aborted near step
  100, because the mean of $c$ (exactly conserved) stays at the round-off of its
  initial size, about $10^{-18}$, while $c$ decays. A net source is still
  rejected (a $10^{-10}$ offset added at $t = 1$ stops the run).
\end{itemize}

\textbf{Results.}
\begin{itemize}
\item \texttt{steady\_transient.i}: $c_t = D\nabla^2 c + k\phi$,
  $-\nabla^2\phi = s c$ on a periodic box, one Fourier mode against its exact
  two-step recurrence; $\phi^0$ the steady solution of $c^0$. Torch PCG: errors
  up to $9\times10^{-15}$ ($c$) and $8\times10^{-17}$ ($\phi$) over 10 steps, on 1, 3
  and 4 ranks; swapping the field order is bit-identical; PETSc agrees to
  $8\times10^{-15}$; MPS float32 $4\times10^{-7}$. Sabotage: no coupling gives
  $3\times10^{-2}$, $\phi^0 = 0$ gives $6\times10^{-3}$.
\item Mean correction: with the $c$ solve alone ($s = 0$), $\overline{c}$ stays
  at $10^{-17}$ over 200 steps at $l_{tol} = 10^{-8}$ and $10^{-6}$ (PCMG), against $-1.5\times10^{-9}$ and
  $1.3\times10^{-8}$ without it. PFHub 1b (no-flux faces) at
  $l_{tol} = 10^{-8}$ keeps the total of $c$ (20101) to $2.5\times10^{-10}$,
  against drifts of $6\times10^{-5}$ (PETSc) and $1.5\times10^{-4}$ (torch PCG)
  without it; new tests at the loose tolerance catch the removal. No
  measurable cost at $256^2$. The existing golds are unchanged within their
  tolerances.
\end{itemize}

\textbf{PFHub 6a} (\texttt{benchmarks/06\_electrostatics}, regression test
\texttt{pfhub\_6a.i}). No new code: $c$ and the steady $\phi$ in one solver,
sources at $t^n$. The face condition $\partial\mu/\partial n = 0$ enters the
update only through the total face flux
$M(\partial B/\partial n - \kappa\,\partial\nabla^2c/\partial n)$, $B = f' + k\phi$, so
the zero split (Neumann faces of $B$ and the biharmonic's Neumann pair) gives the
same $c$ as the benchmark's split to round-off (checked in 1D: difference 0.0).
An inhomogeneous biharmonic face condition was approved and then dropped as
unnecessary (user, 2026-10-26). $F(400)$: 104.00 ($100^2$), 104.01 ($200^2$),
103.82 ($100^2$, solve $\Delta t = 0.05$) on the original schedule of 20 substeps per step
(104.10, 104.11, 103.82 on the smooth schedule of PR7b), against the submissions'
102.9--105.8 (INL MOOSE 103.6, CORAL 104.9, PRISMS-PF 105.8, MMSP 102.9,
MEMPHIS 103.0; FiPy 114.3). Total $c$ conserved to $2\times10^{-13}$ relative.

\status{implemented and verified on branch \texttt{implicit\_fd\_pr5}, tag
\texttt{pr5a-steady-transient}, 2026-10-26. The milestone review (numerics, tests and
parallel, docs, verifier) found no error in the mean correction, the coupling or the 6a
equations (checked against the PFHub notebook). It found: the singular-field mean check
rejecting a neutral decaying source after about 100 steps (fixed with D10, with tests), a
loose-tolerance spec whose error columns could not fail (now gated below $10^{-7}$), the
$t = 0$ free energy of the initial guess of $\phi$ in the PFHub file (left out), and stale
numbers and text in the docs (corrected). Known limit: in single precision the mean correction
is off when $D\,\Delta t/h^2 \gtrsim 3\times10^4$ (2D), where $\lambda$ is below the round-off of
$A\mathbf{1}$. CUDA runs in PR CI.}

\pr{6}{Picard Iteration + Two-Field Coupling (KKS-shaped)}
\textbf{Scope:} iterate the \texttt{ImplicitFDSolver} field loop to
convergence within a time step (design name \texttt{PicardTensorSolver},
not built; whether it wraps \texttt{ImplicitFDSolver} or is an option of it
is open), with \texttt{IterativeTensorSolverInterface}-style knobs (min/max
iterations, tolerances or a Postprocessor, relaxation). Inputs from what is
built: the explicit computes, and so the sources and buffer-valued boundary
values, are evaluated once per time step at $t^n$ (decision~11); a Picard
iteration re-evaluates them, which also removes the boundary-value lag. The
preconditioner rebuild threshold of decision~3 and rebuilds on
solution-dependent coefficients belong here.

\textbf{Verification:} measured Picard convergence on the real coupled case
(R2); cross-validation against \texttt{examples/libtorch\_kks/KKS\_libtorch.i};
per-iteration cost under the reuse policy (R3).

\textbf{Design (user, 2026-10-26).} An option of \texttt{ImplicitFDSolver}, off by
default (today's lagged scheme, unchanged). The coupled step is one nonlinear
system over all fields $U = (u_1, \dots, u_m)$,
$F(U) = A U - b(U) = 0$, with $A$ the sum of each field's
\texttt{LinearFDOperator}s and $b(U)$ the $u^n/dt$ terms plus the explicit
sources evaluated at $U$. Today's field-by-field linear solve with $A$ (PCMG or
torch PCG) is the preconditioner. PETSc SNES drives the iteration, so the method
is a \texttt{petsc\_options} choice: Picard (the default \texttt{newtonls} with the
field solves as preconditioner and the basic line search), Anderson,
NGMRES, and later Newton-Krylov. A plain Picard loop in torch covers the torch-only
backend. The nonlinear driver is chosen independently of the linear solver
(\texttt{nonlinear\_solver = AUTO|PETSC|TORCH}, user, 2026-10-27), so the SNES
methods also run over the torch field solves, e.g.\ on MPS. The SNES vectors are
on the device with a CUDA PETSc and on the host otherwise; host vectors with
device field solves copy every iterate, and the solver warns at the first solve.
Empty \texttt{petsc\_options} lists count as no options.
\begin{itemize}
\item \textbf{D11, ordering.} Every source in an iterate comes from the iterate
  of all fields (Jacobi), so each iterate and the converged result are
  independent of the order of \texttt{buffer}, as D8. Gauss-Seidel was declined.
\item \textbf{D12, convergence.} The single coupled residual norm
  $\|A U - b(U)\|_2$ with the standard SNES relative and absolute tests (rather
  than a per-field scaled norm).
\item \textbf{D13, step.} The full, undamped Picard update by default
  (\texttt{basic} line search, damping 1); damping, line searches and
  Anderson/NGMRES mixing through \texttt{petsc\_options}.
\item \textbf{D14, iteration count.} With the nonlinear iteration on,
  \texttt{getIterations()} (the adaptive time stepper's signal) is the number of
  nonlinear iterations, with the Krylov total printed; with it off, the Krylov
  count as today.
\item \textbf{Coupling to SNES.} Right preconditioning in the
  \texttt{SNESSetPicard} form: the SNES residual is the true $F(U)$ and the
  Jacobian is $A$, applied through each field's existing solve as a block-diagonal
  preconditioner; \texttt{-snes\_mf\_operator} later gives Newton-Krylov with $A$
  as the preconditioner. The existing solve as a nonlinear preconditioner was
  declined: \texttt{nrichardson} would then test the preconditioned residual.
\item \textbf{D15, time level of the explicit computes.} With the nonlinear
  iteration on, the explicit computes are evaluated at $t^{n+1}$, so the converged
  step is fully implicit backward Euler,
  $A U^{n+1} = u^n/dt + f(U^{n+1}, t^{n+1})$, consistent with the boundary
  functions (already at $t^{n+1}$). The lagged scheme (iteration off) keeps
  $f(u^n, t^n)$. A source with an explicit $t$ differs between the two even after
  a single iteration; time-independent sources do not.
\item \textbf{D16, buffer-valued boundary data.} With the nonlinear iteration on,
  the face data of a \texttt{c\_buffer} boundary condition are prepared again after
  every re-run of the explicit computes, so the converged step satisfies
  $c(U^{n+1}, t^{n+1})$ and decision~11's lag is removed. A solution-dependent
  boundary value is then part of the nonlinear coupling, which the block-diagonal
  preconditioner does not see: strong coupling costs iterations or ends in a dt
  cut. Scalar and function boundary data are unchanged.
\end{itemize}
\textbf{Finding (substep 4): the mixing methods need the field solves as their
NPC.} With \texttt{-snes\_type anderson} or \texttt{ngmres} alone, SNES mixes the
raw residuals and never calls the PCSHELL (0 linear iterations): the bicrystal
does not converge in 200 iterations, and on the coupled $c/\phi$ pair an iterate
leaves a mean of $3$--$5\times10^{-13}$ in $\phi$'s right-hand side, above the
singular-field check. With \texttt{-npc\_snes\_type newtonls} (whose linear
solve is the field-solve PC) both converge to round-off. Documented in
\texttt{ImplicitFDSolver.md}; no code change (user, 2026-10-26).

\textbf{Validation: the bicrystal area rate.} \texttt{bicrystal.i} (PFHub 9a's
model) with Anderson, $r_0 = 100$, interface width 30, domain 320, to $t = 20$;
slope of the area fitted over $t \in [5, 20]$, against $-2\pi L\kappa$:
\begin{center}
\begin{tabular}{lrrr}
 & $dt = 0.5$ & $0.25$ & $0.125$ \\
$h = 10$ & & $-3.12\%$ & $-3.35\%$ \\
$h = 5$ & $+1.27\%$ & $+0.74\%$ & $+0.50\%$ \\
$h = 2.5$ & $+2.15\%$ & $+1.62\%$ & $+1.37\%$ \\
$h = 1.25$ & & $+1.83\%$ & $+1.58\%$ \\
\end{tabular}
\end{center}
The spatial error is second order (differences $3.85$, $0.87$, $0.21$ points
between successive $h$) and the time error first order ($0.53$, $0.25$). The
converged deviation is about $+1.4\%$: a property of the model at $r_0 = 100$, not
of the discretisation. The sharp-interface rate holds only for $r \gg$ the interface
width ($r_0/w \approx 3.3$ here), and the deviation grows as the grain shrinks
($+0.9\%$ over $[5,10]$, $+2.1\%$ over $[15,20]$ at $h = 2.5$, $dt = 0.125$). At
$h = 5$ the negative spatial error happens to cancel most of it: that is why the
deviation grew under refinement. Checks: a domain of 480 gives the same slope to
four digits (no periodic-image effect); $r_0 = 200$ ($h = 2.5$, $dt = 0.25$,
domain 640) gives $-0.01\%$.

\textbf{PFHub 6a, fully implicit (substep 5).} \texttt{6a\_implicit\_fd\_nonlinear.i},
a separate input; the lagged input and its results are unchanged (user, 2026-10-27). The
step follows the nonlinear iterations (\texttt{TensorSolveIterationAdaptiveDT}, substeps
$=1$, so a failed solve cuts $dt$). Findings, $100^2$ to $t = 50$, $F$ against a nonlinear
run with $dt \le 0.005$:
\begin{itemize}[leftmargin=*]
\item Plain Picard fails near $t = 41$, where the solve $dt$ passes about $0.19$: the
iterations rise from 3 to the limit of 50 as $dt$ grows (the Picard limit
$dt\,L f'' < 1$). Anderson with the field solves as its NPC reaches $t = 400$ at a solve
$dt$ of 1 with at most 24 iterations.
\item Anderson's iterations grow slowly with $dt$: 4 iterations hold $dt$ at
$1.85\times10^{-3}$, 8--9 at $0.034$, 12 at $0.15$. A step twice as long costs far less
than two steps, so the window was widened from 4--10 to 8--15 (user, 2026-10-27).
\item At the same solve $dt$ the nonlinear error is several times smaller than the lagged
one: $-0.009$ against $+0.039$ at $t = 5$ ($dt = 0.023$), $+0.002$ against $+0.078$ at
$t = 50$ ($dt = 0.25$), both on the lagged input's step schedule (then 20 substeps per step).
\item The iteration window does not see the time error. With 8--15 alone $dt$ reaches
$0.15$ by $t \approx 1$ and $F(5)$ is off by $-0.18$; caps of $0.1$, $0.05$, $0.02$ give
$-0.12$, $-0.060$, $-0.020$ (first order). The input caps at $0.05$ (user,
2026-10-27), a demonstration setting: twice the lagged input's steps at $t = 5$, a fifth of
them at $t = 50$ ($-0.005$). Second order in time is PR7b.
\item An earlier reference ($dt \le 0.025$) carried an error of $-0.009$ at $t = 5$, as large
as the differences it was used to measure there; the numbers above use $dt \le 0.005$.
\item To $t = 400$: 8030 steps, no cuts, 8.9 iterations per step (at most 10),
$F(400) = 103.79$, total $c$ conserved to $3\times10^{-11}$; 1032\,s on one core,
against 18\,s for the lagged input at its larger steps (25\,s on the smooth schedule of PR7b).
\end{itemize}

\textbf{Later (not PR6):} a ``half-Newton'' linearisation that adds each
source's derivative with respect to its own field (a per-cell diagonal from
autograd) to that field's operator, with Newton-Krylov (autograd Jacobian
action) as the fallback. Open there: sources containing derivatives of their
field (Cahn-Hilliard's $\nabla^2 f'(c)$ gives $\nabla^2(f''\cdot)$, not a
diagonal), $f'' < 0$ making the operator indefinite, and the coarse-level
restriction of the diagonal (as D6).

\status{done: tag \texttt{pr6-picard}, milestone review addressed, in PR \#192 (with
PR5 and PR7), 2026-10-27.}

\pr{7}{Port SBM/KKS No-Flux Example to the Implicit FD Path}
\textbf{Scope:} reimplement \texttt{test/tests/kks/KKS\_no\_flux\_bc.i} with
\texttt{TensorNeumannBC} zero-flux faces instead of the \texttt{psi} mask
(\texttt{test/tests/implicit\_fd/kks\_no\_flux.i}).

\textbf{Decisions (user, 2026-10-27).}
\begin{itemize}
  \item \textbf{Geometry.} The solve region of the masked input, the strip
    $-20 < x < 20$, periodic in $y$ over $[-50, 50]$, with sharp zero-flux faces at
    $x = \pm 20$ in place of the diffuse SBM walls (half width $l = 4.2$). Near the walls
    the two differ at order $l$, so the port is a related boundary problem, not a
    reproduction of the masked one. Grid spacing 1 (the masked input's 5 does not resolve
    the interface).
  \item \textbf{Split.} Both parabolas have the curvature $2\rho^2 = 4$ in $c$, so
    $\mu = 4c + g(\eta)$ exactly. $c$: implicit $\partial_t c - 4M\nabla^2 c$ with the
    source $M\nabla^2 g$ (Neumann faces of $g$); $\eta$: implicit
    $\partial_t\eta - L\kappa\nabla^2\eta$ with the source $-L\,\partial F/\partial\eta$ at
    fixed $c$. The PR6 nonlinear solve (Anderson) iterates the sources to the fully implicit
    backward Euler step, in place of the explicit AB3 with 1000 substeps.
  \item \textbf{Verification.} No gated comparison with the masked run and no MOOSE build;
    conservation, energy decrease, convergence, a mirror test and the spectral solver on
    the unmasked periodic problem.
\end{itemize}

\textbf{Results} ($dt = 0.1$, 10 steps, $dx = 1$ unless varied):
\begin{itemize}
  \item total $c$ conserved to $5\times10^{-13}$ (of 1738); the free energy decreases
    every step at every resolution tried.
  \item The strip equals the mirrored initial condition on the doubled periodic domain
    $-20 < x < 60$ (halved integrals) to $5\times10^{-14}$ relative: the test's gold.
  \item $dt$ halved four times: differences of $\int\eta$ and $F$ at $t = 1$ fall by
    2.03--2.11 per halving (first order). $dx$ halved twice: by 4.00--4.07 (second order).
  \item The extrapolated limit ($dt, dx \to 0$) of $\int\eta$ is within
    $7.7\times10^{-2}$ ($3\times10^{-5}$ relative) of the spectral AB3 solver on the doubled
    periodic domain without the mask at $dx = 0.5$, whose own change from $dx = 1$ is
    $7.6\times10^{-2}$. The spectral run at $dx = 0.25$ gives NaN; not investigated.
  \item PETSc multigrid and torch PCG (CPU), 1 and 4 ranks (split along $y$) agree with the gold to
    $10^{-9}$. 26--33 Anderson iterations per step at $dt = 0.1$, 8 at $dt = 0.00625$.
  \item MPS (single precision): the nonlinear residual stalls at $10^{-5}$ of its initial
    value in the second step and the iteration does not converge at
    \texttt{nl\_rel\_tol} $= 10^{-5}$ or $5\times10^{-5}$; no MPS test. Open.
\end{itemize}

\textbf{Example} (\texttt{examples/implicit\_kks/KKS\_no\_flux\_implicit.i}): to $t = 100$
with \texttt{TensorSolveIterationAdaptiveDT}, 682 steps, total $c$ conserved to
$4\times10^{-12}$, energy non-increasing (88.43 to 42.02). As the grain relaxes the initial
residual of a step falls to $10^{-2}$ and \texttt{nl\_rel\_tol} $= 10^{-10}$ of it is below
the residual's round-off ($10^{-12}$), so with the relative test alone the iteration stalled
and the step was cut to \texttt{dtmin} at $t = 40.8$. The example adds
\texttt{nl\_abs\_tol} $= 10^{-9}$; the test input keeps the relative test only.

\status{done: tag \texttt{pr7-kks-no-flux}, in PR \#192, 2026-10-27. The milestone review
(numerics with an independent Newton backward-Euler reference, docs and parallel, mutation
testing, verifier) found no error in the model or the split. It found a conservation gate
that could not fail and a CUDA test skipped behind a CPU-only prerequisite (both fixed), and
a redundant ghost exchange of the nonlinear iterate (removed; the one after the linear solve
is an open question). CUDA runs in CI.}

\pr{7b}{Second-Order Time Integration (BDF2)}
\textbf{Scope (user, 2026-10-27):} BDF2 for the transient fields of
\texttt{ImplicitFDSolver}, $(3u^{n+1} - 4u^n + u^{n-1})/(2\,dt) + A u^{n+1} = b(U^{n+1})$,
with the variable-step coefficients, since the adaptive stepper changes $dt$.
With the nonlinear iteration of PR6 it is fully implicit to second order. Backward
Euler's first-order time error is what bounds $dt$ in the PFHub 6a study of PR6, not
the nonlinear convergence (Anderson converges in 8--15 iterations at steps that
carry a visible time error), so this, and not more iterations, is how to take
larger steps at a fixed accuracy.

\textbf{Design questions (to decide):} where $u^{n-1}$ and the previous $dt$ live
(\texttt{TimeDerivativeOperator} with its own history, or the solver's stored $u^n$
of PR6); the first step (backward Euler, or a smaller start); a step cut, which must
restore the history; the coefficient change of the operator, $3/(2dt)$ in place of
$1/dt$, which changes the stencil and so rebuilds the preconditioner when $dt$ or the
step ratio changes (decision~3); the lagged scheme of PR3--PR5a (whether BDF2 is
offered there, where the explicit terms are first order anyway); and an error
estimate (BDF2 against backward Euler) for an accuracy-driven $dt$, as an
alternative to the iteration-driven stepper.

\textbf{Verification:} order 2 in $dt$ on the closed-form coupled recurrence of PR6
and on the scalar Newton reference (the BDF2 recurrence); order 2 of $F(t)$ in the
6a study; mass conservation to round-off.

\textbf{Substeps removed (first PR7b step, user, 2026-10-28).} \texttt{ImplicitFDSolver}
no longer takes \texttt{substeps}: one solve per time step, so a failed solve can always be
retried with a smaller $dt$, and a step history holds the previous step. The inputs that
used substeps (the Crank tests, PFHub 1a/1b/6a and their regression tests) now take the
smooth equivalent schedule: the first $dt$ and \texttt{dtmax} divided by the former substep
count $s$, the growth factor $1.1^{1/s}$ per step (the same growth over $s$ steps), and
output every $s$ steps through a new \texttt{time\_step\_interval} of the tensor outputs
(with \texttt{FINAL} for the last step). The Crank results are unchanged to $2\times10^{-7}$
relative (the same steps). The 1b and 6a regression golds were regenerated (the solve steps
differ. 1b shares no output times with its old gold: $F$ interpolated to the old times differs
by up to $2\times10^{-4}$ relative. 6a shares $t = 1$, where $F$ differs by $6\times10^{-7}$
relative, and about $2\times10^{-6}$ interpolated to the old times after the first step). PFHub 6a on the smooth schedule: $F(400) = 104.10$ ($100^2$), 104.11
($200^2$), 103.82 (cap 0.05); against the cap of 0.05, caps 0.5, 0.25, 0.1 differ by
0.280, 0.083, 0.019.

\textbf{BDF2 as built (decisions by the user, 2026-10-28).} \texttt{TimeDerivativeOperator}
\texttt{order = 2}: MOOSE's variable-step weights ($s = dt + dt_\text{old}$: $w_0 = 1 + dt/s$,
$w_1 = -s/dt_\text{old}$, $w_2 = dt^2/(dt_\text{old} s)$), with a centre stencil entry of
$w_0/dt$ and a contribution of $-(w_1 u^n + w_2 u^{n-1})/dt$. The first step is backward Euler.
The operator keeps $u^n$, $u^{n-1}$ and $dt_\text{old}$ itself (the buffer states have no
$u^0$). They are keyed by the start time of the step, so a retried step keeps them. It is
allowed with the lagged scheme and documented there as first order overall when there are
explicit terms. The solver is unchanged: a $w_0$ change is a stencil change, so the existing
rebuild covers it (decision~3). Verification:
\begin{itemize}[leftmargin=*]
\item A periodic mode over 8 steps with ratios 0.5 to 2, against the exact BDF2 recurrence
(\texttt{bdf2\_mode.py}): $3\times10^{-8}$ (PETSc), $4\times10^{-10}$ (torch) relative. At
fixed steps to $t = 0.1$ the error against the exact decay falls at order 2.24, 2.13, 2.06
(backward Euler 1.06, 1.04, 1.02); manual runs of \texttt{bdf2\_mode.i} with
\texttt{TensorComputes/Solve/dudt/order} 1 or 2,
\texttt{Executioner/TimeStepper/type=ConstantDT}, \texttt{Executioner/TimeStepper/dt} $= 0.1/m$
and \texttt{Executioner/num\_steps} $= m$ for $m = 10, 20, 40, 80$,
\texttt{TensorSolver/l\_tol=1e-13} and \texttt{--allow-unused}.
\item The bicrystal with the nonlinear iteration on varying steps, against a Newton BDF2
reference: $4\times10^{-14}$.
\item A step cut with the history in place equals a run on the accepted times from the start.
\item Mutants caught: a retry that advances the history, constant-step weights, and a flipped
$w_2$. A stencil without $w_0$ survives (it changes only the multigrid preconditioner, not the
solution).
\item PFHub 6a, fully implicit, $100^2$. $F$ against BDF2 at $dtmax = 0.005$, for dtmax 0.1, 0.05
and 0.02: backward Euler $-0.130$, $-0.067$, $-0.027$ at $t = 5$; BDF2 $-3.8\times10^{-4}$,
$-1.1\times10^{-4}$, $-1.9\times10^{-5}$ (at $t = 1$: $-2.1\times10^{-4}$, $-6.0\times10^{-5}$,
$-8.1\times10^{-6}$, order 1.8--2.2). Same steps, no cuts, and wall time within 10\%. Manual
runs of \texttt{6a\_implicit\_fd\_nonlinear.i} with \texttt{n=100 end\_time=5},
\texttt{Executioner/dtmax} and \texttt{TensorComputes/Solve/dcdt/order}.
\item Restart: the history is not restartable data, so a recovered run would take backward
Euler on its first step. Recover does not restore the tensor buffers either (a known issue).
\end{itemize}

\status{done (2026-10-28), tag pr7b-bdf2, commits ``(PR7b)'' in \#192.}

\pr{9}{Smoothed Boundary Method on the Implicit FD Path}
\textbf{Scope:} the smoothed boundary method (Yu, Chen and Thornton, arXiv:1107.5341) on
\texttt{ImplicitFDSolver}, with the KKS case of \texttt{test/tests/kks/KKS\_no\_flux\_bc.i}
(no flux, $90^\circ$ contact angle) as the target. Branch \texttt{implicit\_fd\_sbm}, on top of
\#192; commits ``(PR9)''.

\textbf{Equations.} For no flux (paper Eq.~4--6 with $B_N = 0$, Eq.~22 at $\theta = 90^\circ$),
multiplied through by the domain parameter $\psi$:
\begin{align*}
\psi\,\partial_t c &= \nabla\cdot(\psi M\nabla\mu), &
\psi\,\partial_t \eta &= \nabla\cdot(\psi L\kappa\nabla\eta) - \psi L\,\partial F/\partial\eta.
\end{align*}
With $\mu = 4c + g(\eta)$ (PR7): implicit $\psi\partial_t c - \nabla\cdot(4\psi M\nabla c)$, source
$\nabla\cdot(\psi M\nabla g)$; implicit $\psi\partial_t\eta - \nabla\cdot(\psi L\kappa\nabla\eta)$,
source $-\psi L\,\partial F/\partial\eta$. The operators are symmetric and need no division by
$\psi$. $\sum\psi c$ is conserved to round-off. A nonzero boundary flux, a contact angle
other than $90^\circ$ (both $|\nabla\psi|$ sources), Dirichlet boundaries (Eq.~7) and the
Cahn-Hilliard form with gradient energy (Eq.~23, whose fourth-order operator is not a
symmetric product of $\psi$-weighted operators) are out of scope.

\textbf{Decisions (user, 2026-10-28).}
\begin{itemize}[leftmargin=*]
\item The $\psi$-multiplied form, with $c$ (not $\psi c$) as the unknown.
\item A new operator \texttt{CoefficientTimeDerivativeOperator}, $a(\mathbf{x})\,\partial_t u$
with a \texttt{coefficient} buffer and the backward Euler and BDF2 weights of
\texttt{TimeDerivativeOperator}; its stencil centre is the per-cell field $a w_0/dt$, restricted
to the coarse levels as the \texttt{FiniteDifferenceDivGrad} coefficient is. It is reused for
grand-potential models ($\chi\,\partial_t\mu$).
\item $\psi$-weighted sources are ordinary sources formed in the input.
\item The $\psi$-weighted diffusion uses \texttt{FiniteDifferenceDivGrad} with
\texttt{face\_mean = arithmetic}: a $\psi = 0$ cell next to the domain has a positive face
coefficient and no time-derivative weight, so it acts as a ghost cell carrying the no-flux
condition.
\item \textbf{Inactive cells.} A cell with a zero diagonal and a zero right-hand side (so
$\psi = 0$ on the cell and on all its faces) is inactive: $u$ is held at $u^n$, an identity row
in the operator action, the Jacobi diagonal, the PETSc levels and the residual. A zero diagonal
with a nonzero right-hand side stays an error, and so does a zero-diagonal cell whose row is not
zero ($A$ applied to the inactive indicator must vanish). A coarse cell is inactive only if all
its children are.
\item The threshold for small $\psi$ stays in the input, as in the spectral input: a $\psi$
profile that reaches 0 (the cosine ramp \texttt{psi\_func}), or a \texttt{ParsedCompute} that
sets $\psi = 0$ below a cutoff.
\item From the spectral \texttt{ReciprocalMatDiffusion} and \texttt{ReciprocalAllenCahn}: the
profile, the active set $\psi > 0$, and the held state outside the domain are kept. The FFT
derivatives, the separate $\nabla\psi/\psi$ term (contained in the multiplied form) and
\texttt{always\_update\_psi} (the coefficient version check covers it) are not.
\end{itemize}

\textbf{Constraint: $\psi$ will change in time} (user, 2026-10-28; e.g.\ a KKS-like model
recast as SBM, with thermodynamics coupled at the interfaces only, and Gibbs energy minimisers).
Nothing about $\psi$ is fixed at setup: the inactive set, the stencils and the coarse
coefficients follow every change of the coefficient buffers, and a cell that becomes active
starts from its held value. Open for then (a physics decision): with $\partial_t\psi \ne 0$,
$\psi\,\partial_t c$ does not conserve $\int\psi c$; the conservative form is $\partial_t(\psi c)$.

\textbf{Steady singular fields} (e.g.\ Poisson with SBM no-flux faces): the nullspace is the
constants on each connected active region, not the whole grid. Detected and an error until
needed.

\textbf{Verification:}
\begin{itemize}[leftmargin=*]
\item \texttt{CoefficientTimeDerivativeOperator} with $a = 1$ equals
\texttt{TimeDerivativeOperator} (backward Euler and BDF2); a spatially varying $a > 0$ against
an independent sparse solve of the same discretisation.
\item Inactive cells: held at $u^n$ exactly, the active cells' solution unchanged by the
inactive block, both backends, 1--4 ranks; the error paths.
\item A 1D slab with $\psi$ walls against the sharp zero-flux solution, converging as the wall
width $l \to 0$.
\item The KKS case against the spectral SBM input (same $\psi$, same $l$): $\sum\psi c$
conserved to round-off, the free energy non-increasing, and $\eta$, $c$ close to the spectral
run, limited by the spatial discretisation.
\end{itemize}

\textbf{Record (2026-10-28).}
\begin{itemize}[leftmargin=*]
\item \texttt{coefficient\_dt.i}: $a = 1$ reproduces the BDF2 mode gold; with
$a = 1 + 0.5\sin 2\pi x\sin 2\pi y$ both backends match a scipy sparse solve
(\texttt{coefficient\_dt.py}) to $1.4\times10^{-8}$, and $\int a u$ is conserved.
\item \texttt{sbm\_inactive.i}: the inactive cells stay at $u^n$ exactly; the active cells
match scipy on the active block (\texttt{sbm\_inactive.py}) to $3\times10^{-9}$ (PETSc, 4 ranks),
$2\times10^{-12}$ (torch) and $2\times10^{-14}$ (the tested torch nonlinear iteration;
$2.5\times10^{-11}$ with the PETSc one); in single precision, $7\times10^{-7}$ (CPU, MPS).
\item \texttt{sbm\_slab.i}: the error against the sharp zero-flux solution at $t = 0.1$ is
$3.7\times10^{-4}$, $1.05\times10^{-4}$, $2.5\times10^{-5}$, $6.1\times10^{-6}$ and
$1.8\times10^{-6}$ for $l = 0.08$ to $0.005$ (halving, 4 cells per $l$): order 2 in $l$.
\item \texttt{kks\_sbm.i}: $\int\psi c$ conserved to $10^{-12}$, the free energy decreasing,
both backends and 1 and 4 ranks within $10^{-13}$. At $t = 1$, $\int\psi\eta$ is 2328.73
(time-converged, $dt = 0.01$) against 2329.46 for the spectral AB3 run
(\texttt{KKS\_no\_flux\_bc.i}, the same grid and $\psi$), a gap of 0.73 in a change of 129; at
$dx = 0.5$ (time-converged, $dt = 0.01$ and $0.005$ within $10^{-3}$), 2329.355 against 2329.539
(gap 0.18).
\item \textbf{Identity-row scale (preconditioner only).} With the largest diagonal on the
inactive identity rows, PCMG took about 100 iterations per solve on \texttt{kks\_sbm.i}; with the
smallest active diagonal it takes 12--13, as with $\psi = 1$ (12) or a $\psi$ floor without
inactive cells. The inactive rows and columns are decoupled from the active block, so the
solution is unchanged (all golds pass unchanged).
\item \textbf{Review (2026-10-29).} Aspect review, mutation testing (17 mutations, 9 killed, two of them by
tests added in the review: a PCMG iteration cap of 15 on \texttt{kks\_sbm.i}, serial (14--16 on 4
ranks), and the within-step coefficient change error. Of the 8 survivors, 7 change only the
PCMG preconditioner or the speed of the torch PCG solve; the eighth disables the error for a
zero-diagonal cell with nonzero off-diagonal entries, which no test reaches)
and a fix verification. Fixed: the single-precision identity scale (DBL\_MAX in float32), the
quoted figures.
\end{itemize}

\status{done (2026-10-29, tag \texttt{pr9-sbm}); new PR from branch \texttt{implicit\_fd\_sbm},
on top of \#192.}

\pr{8 (deferred)}{Monolithic Block-Matrix Groundwork}
\textbf{Scope:} not scheduled. Revisit only if a model needs exact
(non-Picard) multi-field coupling; it would use field-split / block
\texttt{Vec}s (Section~\ref{sec:soa}). The \texttt{target\_field} tag
proposed for it in PR0 was not built.

\status{deferred until requested.}

\subsection*{Outside the core PRs: PFHub 2a with the implicit FD solver}
On branch \texttt{pfhub\_2a\_implicit\_fd} (on top of PR4); its PR goes up
together with 2b, which needs PR5's no-flux biharmonic (user, 2026-10-25). The existing 2a inputs started $c$
with $\epsilon = 0.01$; the PFHub specification uses $0.05$, corrected on
that branch. Results posted with $0.01$ (\texttt{results.2a.1.json}) and the
notebook's \texttt{2a\_secant\_out.csv} need regenerating.

\section{Status Tracking}

\begin{longtable}{>{\bfseries}p{0.08\textwidth} p{0.34\textwidth} p{0.26\textwidth} p{0.22\textwidth}}
\toprule
PR & Deliverable & Status & Verification \\
\midrule
\endhead
PR0 & LinearFDOperator mixin + Laplacian retrofit & merged (\texttt{2580075}, PR \#185) & stencil vs.\ conv-kernel unit tests \\
PR1 & BC TensorComputes, ghost fill & merged (PR \#186) & real-space explicit tests; dual-path unit test (rewritten in PR4) \\
PR1a & Torch-native vs.\ PETSc matrix-free benchmark & spike, reviewed; not merged (tag \texttt{pr1b-spikes}) & CPU+CUDA logs and tables (PR1a) \\
PR1b & Multigrid spikes $\to$ hybrid design & spike, reviewed; not merged (tag \texttt{pr1b-spikes}) & chain table and logs (Section~\ref{sec:pr1b}) \\
PR2 & DMDA/decomposition compatibility & spike, not merged; superseded by \texttt{PetscFDHierarchy} & 1--8 ranks; ownership checked in production \\
PR2b & Multi-rank PCMG & spike done (\texttt{b3dcd72}) & PETSc-owned hierarchy + cell-centred transfer adopted (user, 2026-10-18) \\
PR2c & Non-SPD + torch PCSHELL smoother + fusion & spike done (\texttt{40f6977}, \texttt{65109ae}) & GS robust for upwind; central Pe $\ge1$ deferred; CUDA traced 2.4--3.2$\times$ eager \\
PR3 & ImplicitFDSolver, periodic, matrix-free & merged (PR \#188) & MMS and BE orders 2/4; CPU 1--4 ranks, MPS, H200 \\
PR3b & Biharmonic, PFHub 1a input, smoother choice (Chebyshev default), no divergence test, preallocation, scaling script & merged (PR \#189) & biharmonic Fourier mode 2D/3D, both backends \\
PR4 & BCs in ImplicitFDSolver, face-varying values & implemented, pre-push reviewed, findings addressed; CPU serial and multi-rank pass; PR \#191 open, to merge before \#192 (tag \texttt{pr4-implicit-bcs}) & MMS vs.\ direct solve, both backends; PCMG iterations flat in $n$ \\
PR4a & AMG mixed mode, contingent & proposed & --- \\
PR5 & Variable coefficients, no-flux biharmonic, PFHub 1b & done (tag pr5-variable-coefficient), in PR \#192 & DivGrad MMS (2D, 3D, coefficient jump), capacitor, biharmonic faces MMS; CPU 1--4 ranks \\
PR5a & Steady + transient coupling & done (tag pr5a-steady-transient), in PR \#192 & sources at $t^n$ (order-independent); exact mean of the solution; coupled test against the exact recurrence; PFHub 6a within the submissions' spread \\
PR6 & Picard, two-field coupling & done (tag pr6-picard), in PR \#192 & coupled BE recurrence, time source, bicrystal Newton reference; CPU 1--4 ranks, MPS \\
PR7 & SBM/KKS no-flux port & done (tag pr7-kks-no-flux), in PR \#192 & mirror gold, conservation, energy decrease, $dt$/$dx$ orders, spectral limit; CPU 1--4 ranks \\
PR7b & Substeps removed, output interval, BDF2 time integration & done (tag pr7b-bdf2), in PR \#192 & exact variable-step recurrence, order 2, Newton BDF2 reference, cutback history, 6a \\
PR8 & Monolithic groundwork & deferred & --- \\
PR9 & Smoothed boundary method, \texttt{CoefficientTimeDerivativeOperator}, inactive cells & done (tag pr9-sbm), new PR on top of \#192 & coefficient $dt$ and inactive block vs.\ scipy, slab order 2 in $l$, KKS vs.\ spectral SBM; CPU 1--4 ranks \\
--- & PFHub 2a, implicit FD & branch \texttt{pfhub\_2a\_implicit\_fd}; PR together with 2b & --- \\
\bottomrule
\end{longtable}

\section{Open Questions}

\begin{enumerate}[leftmargin=*]
\item Preconditioner rebuild threshold of decision~3: its value and whether
it is user-facing (PR6).
\item DMDA stencil width when operators need different ghost widths: one DM
sized to the widest stencil (as built); revisit if profiling shows waste.
\item Whether PCMG with point Jacobi suffices on the KKS coefficients (R7,
PR4a).
\item Non-SPD operators (advection, central differencing at high Pe):
deferred until a target physics needs implicit advection (PR2c).
\item A Newton direction for the nonlinear iteration (PR6 plan): each source's derivative
with respect to its own field as a per-cell diagonal of that field's operator, with
Jacobian-free Newton--Krylov for cross-field derivatives. Open: sources with derivatives of
$u$ (not diagonal), $f'' < 0$ (indefinite), and the coarse-level restriction.
\item A per-cell reaction operator $k(\mathbf{x})\,u$ on the solved field, with $k$ from the
other fields, for grand-potential models (user, 2026-10-28). The nonlinear iteration already
refreshes such a coefficient at every iterate.
\item An accuracy-driven $dt$ from a BDF2 against backward Euler estimate (deferred from PR7b).
\item Restart: \texttt{--recover} restores no tensor buffers, and the BDF2 history is not
restartable data.
\item Work on steps that write no output: \texttt{TensorProblem::executeTensorOutputs} runs at
every \texttt{TIMESTEP\_END}, so the postprocess computes, the CPU copies and the mapping to
MOOSE variables run at every time step, not every output step. With substeps removed (PR7b)
that is every solve. Whether to skip the copies and the mapping on such steps is open.
\item The ghost exchange of each field after its linear solve
(\texttt{ImplicitFDSolver::substep}) is likely redundant: every compute that reads a
field's ghosts exchanges them itself (\texttt{runComputeWithGhosts}), and removing the
matching exchange from the nonlinear iteration left every \texttt{implicit\_fd} output
byte-identical on 1--4 ranks (PR7 review). Remove it once the same check (serial, MPI, the
face and \texttt{c\_buffer} tests) passes without it; add communication back only where a
reader needs it.
\end{enumerate}

\textbf{Resolved:} input syntax for several fields (one
\texttt{ImplicitFDSolver} with \texttt{buffer} and \texttt{;}-separated
\texttt{operators}; PETSc options apply to every field under its own prefix
\texttt{<field>\_}, not per-field sub-blocks); whether to keep the DMDA
(yes: PETSc owns the hierarchy, Section~\ref{sec:hybrid}); PETSc or
torch-only (PETSc on CPU and CUDA, 2026-10-18); MPS (torch PCG with a Jacobi
preconditioner; the spike torch smoothers are not used); coarse BC rows
(rediscretised, decision~10).

\end{document}
